% -*- coding: utf-8 -*-
% !TEX program = xelatex
%%
%%  NKU Master's Thesis
%%  Multi-Agent Knowledge-Graph-Guided Reasoning for Reliable Log-Based Root Cause Analysis
%%
%%  Compile with XeLaTeX (recommended):
%%    xelatex main && bibtex8 main && xelatex main && xelatex main
%%
\documentclass[12pt,openany]{book}
\usepackage{amsmath,amssymb}
\usepackage{pifont}
\usepackage{ifxetex}
\ifxetex
  \usepackage[bookmarksnumbered]{hyperref}
\else
  \usepackage[unicode,bookmarksnumbered]{hyperref}
\fi

%% English mode NKU thesis
\usepackage[emptydoublepage,English]{NKThesis}

%% Bibliography: biblatex with nkthesis style
\usepackage[backend=bibtex8, defernumbers=true, sorting=nty, style=nkthesis]{biblatex}

\hypersetup{colorlinks=true,
            pdfborder=0 0 1,
            citecolor=black,
            linkcolor=black}

\usepackage{tikz}
\usepackage{pgfplots}
\pgfplotsset{compat=1.18}
\usetikzlibrary{shapes.geometric, arrows.meta, positioning, calc}

\usepackage{booktabs}
\usepackage{array}
\usepackage{multirow}
\usepackage{longtable}
\usepackage{algorithm}
\usepackage{algpseudocode}
\usepackage{listings}
\usepackage{float}
\usepackage{pdfpages}

\lstset{
  basicstyle=\ttfamily\small,
  breaklines=true,
  frame=single,
  numbers=left,
  numberstyle=\tiny,
  keywordstyle=\color{blue},
  commentstyle=\color{green!60!black},
  stringstyle=\color{red},
  showstringspaces=false
}

\addbibresource{references.bib}
\DeclareBibliographyCategory{cited}
\AtEveryCitekey{\addtocategory{cited}{\thefield{entrykey}}}

\includeonly{
  abstract,
  chapters/01_introduction,
  chapters/02_related_work,
  chapters/03_system_design,
  chapters/03b_methodology,
  chapters/04_implementation,
  chapters/05_evaluation,
  chapters/06_discussion,
  chapters/07_conclusion,
  appendices,
  references,
  acknowledgements,
  resume
}

%% Custom commands
\newcommand{\systemname}{AetherLog 2.0}
\newcommand{\fone}{F\textsubscript{1}}

%% Theorem environments
\newtheorem{Theorem}{\hskip 2em Theorem}[chapter]
\newtheorem{Lemma}[Theorem]{\hskip 2em Lemma}
\newtheorem{Corollary}[Theorem]{\hskip 2em Corollary}
\newtheorem{Proposition}[Theorem]{\hskip 2em Proposition}
\newtheorem{Definition}[Theorem]{\hskip 2em Definition}
\newtheorem{Example}[Theorem]{\hskip 2em Example}
\newtheorem{Remark}[Theorem]{\hskip 2em Remark}

\begin{document}

%% ---------------------------------------------------------------
%%  Thesis metadata
%% ---------------------------------------------------------------
\NKTsetup{%
  论文题目(中文) = 基于多智能体知识图谱引导推理的可靠日志根因分析,
  副标题        = ,
  论文题目(英文) = Multi-Agent Knowledge-Graph-Guided Reasoning
                   for Reliable Log-Based Root Cause Analysis,
  论文作者       = ,
  学号           = ,
  指导教师       = ,
  指导教师职称   = ,
  申请学位       = 学历硕士,
  培养单位       = ,
  学科专业       = 计算机科学与技术,
  研究方向       = 人工智能与系统软件,
  答辩委员会主席  = ,
  评阅人         = ,
  中图分类号     = TP311,
  UDC            = 004,
  学校代码       = 10055,
  密级           = 公开,
  保密期限       = ,
  审批表编号     = ,
  批准日期       = ,
  论文完成时间   = 二〇二六年一月,
  答辩日期       = ,
  签字日期       = 2026\quad 年\quad 1 \quad 月,
  论文类别       = 学历硕士,
  院/系/所       = ,
  专业           = ,
  联系电话       = ,
  Email          = ,
  通讯地址(邮编) = ,
  备注           =
}

% -*- coding: utf-8 -*-

\begin{zhaiyao}

现代分布式系统产生大量日志数据，这些数据对于诊断故障和执行根因分析（RCA）至关重要。
基于规则和统计的方法难以扩展到异构云原生环境，且通常缺乏可操作的解释。
大型语言模型（LLM）为日志理解和解释生成提供了新的接口；
然而，单一LLM的RCA存在幻觉、上下文有限以及缺乏验证等问题。

本文提出\systemname{}，一种多智能体、知识图谱引导的RCA系统，
通过以下方式提高可靠性：
（i）多个LLM智能体之间的专业化分工；
（ii）具有显式评判机制的结构化辩论协议。
系统由六个智能体组成：日志解析器、知识图谱检索智能体、
三个RCA推理智能体（日志聚焦、KG聚焦和混合型）
以及一个评判智能体，该智能体根据证据质量、推理强度、置信度校准、
完整性和一致性对竞争假设进行评分。

实验在三个数据集上进行评估：
LogHub的Hadoop1（55个标记应用，4类故障）、
LogKG的CMCC（93个工业OpenStack故障案例，7类）
以及HDFS\_v1（来自575,061个HDFS块跟踪的200个平衡样本，二进制异常检测）。
在Hadoop1上，\systemname{}粗粒度宏观F1比单智能体基线提高了13.3个百分点（39.1\% vs 25.8\%）。
在CMCC上，多智能体系统达到61.3\%的准确率，
而单智能体和RAG基线均崩溃为"未知"或多数类预测。
在HDFS\_v1上，\systemname{}实现了69.5\%的准确率和69.4\%的平衡宏观F1。

实验结果支持多智能体辩论通过交叉验证假设、
减少多数类崩溃并产生基于证据的解释来提高事实性和鲁棒性的假设。

\end{zhaiyao}

\begin{guanjianci}
根因分析；日志分析；大型语言模型；多智能体系统；知识图谱；辩论协议；AIOps
\end{guanjianci}


\begin{abstract}

Modern distributed systems generate massive volumes of logs that are essential for diagnosing
failures and performing root cause analysis (RCA). While rule-based and statistical approaches
can detect known patterns, they struggle to scale to heterogeneous cloud-native environments
and often lack actionable explanations. Recent Large Language Models (LLMs) provide a new
interface for log understanding and explanation generation; however, single-LLM RCA suffers
from hallucinations, limited context, and the absence of verification.

This thesis presents \systemname{}, a multi-agent, knowledge-graph-guided RCA system that
improves reliability through (i)~specialization across multiple LLM agents and (ii)~a
structured debate protocol with an explicit judge. The system consists of six agents:
a Log Parser for structured extraction, a Knowledge Graph (KG) Retrieval Agent for
historical grounding, three RCA Reasoners (log-focused, KG-focused, and hybrid), and a
Judge Agent that scores competing hypotheses on evidence quality, reasoning strength,
confidence calibration, completeness, and consistency. The knowledge layer is implemented
in Neo4j and stores incidents, entities, and root-cause relations to support retrieval and
recurrence analysis.

We evaluate \systemname{} on three datasets: Hadoop1 from LogHub (55 labeled applications,
4 fault classes), CMCC from LogKG (93 industrial OpenStack failure cases, 7 classes), and
HDFS\_v1 (200 balanced samples from 575,061 HDFS block traces, binary anomaly detection).
On Hadoop1, \systemname{} improves coarse macro-\fone{} by +13.3 percentage points over a
single-agent baseline (39.1\% vs 25.8\%) and detects minority faults such as disk\_full
(36.4\% \fone{} vs 0\%). On CMCC, the multi-agent system achieves 61.3\% accuracy while
the single-agent and RAG baselines collapse to ``unknown'' or majority-class predictions,
providing strong evidence for cross-domain generalizability on real industrial logs.
On HDFS\_v1, \systemname{} achieves 69.5\% accuracy and a balanced macro-\fone{} of 69.4\%,
outperforming single-agent (65.0\%) and RAG (63.0\%) baselines.

Overall, the results support the hypothesis that multi-agent debate improves factuality
and robustness by cross-validating hypotheses, reducing majority-class collapse, and
producing evidence-grounded explanations. We discuss limitations such as label ambiguity,
dataset-specific KG coverage, and computational overhead, and we outline future work
including larger-scale ablations, stronger retrieval, and additional ML/DL baselines.

\end{abstract}

\begin{keywords}
Root Cause Analysis, Log Analysis, Large Language Models, Multi-Agent Systems,
Knowledge Graphs, Debate, AIOps
\end{keywords}

\tableofcontents

% -*- coding: utf-8 -*-

\chapter{Introduction}

\section{The Role of Logs in Modern Operations}

Logs are the most pervasive and continuously available telemetry in modern software systems.
Unlike metrics, which aggregate numerical measurements, or distributed traces, which require
explicit instrumentation at the call level, logs are emitted by virtually every system component
as a natural by-product of execution~\cite{he2020loghub}. They record timestamped events,
exception messages, state transitions, and resource identifiers in semi-structured natural
language, making them the primary artifact that on-call engineers consult when a failure occurs.
In cloud-native architectures that span microservices, container orchestration platforms, and
distributed storage systems, a single production incident may leave traces across dozens of
components simultaneously, each producing its own log stream. Correlating and interpreting these
streams within the tight time windows demanded by service-level objectives (SLOs) is a task that
has grown in difficulty commensurate with the growth of system scale and
complexity~\cite{zhang2024cloudRCAsurvey}.

The practical importance of log analysis has motivated a substantial body of research in automated
log parsing, anomaly detection, and failure diagnosis. Log parsers such as Drain convert raw,
free-form log messages into structured event templates by masking variable
tokens~\cite{he2017drain,zhu2019tools}. These templates serve as the input to downstream analytical
pipelines, enabling quantitative comparison of log behaviors across time windows. However, log
analysis and root cause analysis (RCA) are distinct tasks with different objectives: while log
analysis broadly encompasses anomaly detection, classification, and summarization, RCA is
specifically concerned with identifying the underlying causal fault that produced an observable
incident and providing an actionable explanation for operators. This thesis is concerned
exclusively with the RCA task.

\section{Root Cause Analysis: Definition and Challenges}

Root cause analysis is the process of tracing an observed failure back to its primary causal
factor, as distinct from its immediate symptoms. In a microservice system, for example, a
database connection timeout may surface as an HTTP 500 error in an upstream API; the symptom is
the HTTP error, but the root cause is the database unavailability. Correct identification of the
root cause is prerequisite to effective remediation: treating the symptom (e.g., restarting the
upstream service) without addressing the root cause (e.g., the exhausted database connection pool)
will not prevent recurrence~\cite{meng2020localizing,wu2020microrca}.

The RCA task in production environments faces several well-documented challenges that are
difficult to address simultaneously. First, production systems emit logs at a rate that makes
manual inspection infeasible. A single Hadoop cluster, for example, can emit hundreds of thousands
of log lines per job execution, as evidenced by the Hadoop1 benchmark used in this
thesis~\cite{he2020loghub}. Engineers cannot review this volume manually during an active incident.
Second, failures in distributed systems rarely confine themselves to a single component. Fault
propagation across service dependencies means that the failure signal in the logs of the faulty
component may be temporally and spatially separated from the logs of the component that originally
triggered the cascade~\cite{wang2018cloudranger,chen2014causeinfer}. Third, the heterogeneity
of log formats across different software stacks, versions, and configurations means that pattern
matching rules and templates crafted for one system are rarely transferable to another without
substantial manual effort~\cite{zhang2019logrobust}. Fourth, logs encode domain-specific
terminology and operational conventions that require contextual knowledge to interpret correctly.
An error string like ``BlockInfo not found in volumeMap'' is diagnostic only to an engineer who
knows the semantics of HDFS block management. Capturing and operationalizing this domain knowledge
at scale is a fundamental bottleneck in automated RCA.

Traditional automated RCA approaches address these challenges through two broad strategies.
The first class of approaches employs statistical correlation and causal graph methods: they
construct dependency graphs between services or metrics and identify components whose
behavior changes are most causally associated with the observed
anomaly~\cite{wang2018cloudranger,meng2020localizing,chen2014causeinfer}. While effective in
controlled environments where component relationships are well-modeled, these methods require
significant upfront engineering effort to build and maintain the causal model, and they are
brittle when system topology changes frequently. The second class of approaches applies supervised
machine learning to structured log representations: log templates are extracted, converted into
feature vectors or event sequences, and fed into classifiers or sequence models trained to
recognize fault patterns~\cite{xu2009consolelogs}. These methods can achieve strong performance
on benchmark datasets, but they depend on labeled training data from the target system, are
sensitive to log statement drift caused by software updates~\cite{zhang2019logrobust}, and
produce predictions without interpretable explanations---a significant drawback when the goal is
to support operator decision-making rather than merely raise an alert.

\section{LLM-Based Attempts at RCA: Capabilities and Limitations}

The emergence of large language models (LLMs) with strong instruction-following and natural
language generation capabilities~\cite{brown2020languagemodelsfewshotlearners,ouyang2022traininglanguagemodelsfollow}
has motivated a new generation of RCA approaches that treat the diagnosis task as a prompted
reasoning problem. The central appeal is clear: an LLM can read heterogeneous log text without
requiring a domain-specific parser, generate a human-readable explanation of the likely root cause,
and suggest remediation steps---all in a zero-shot or few-shot setting that requires no task-specific
training data.

Several recent systems demonstrate the viability of this direction. RCACopilot, developed at
Microsoft and evaluated on a year's worth of production cloud incidents, employs an LLM to
aggregate diagnostic information collected by per-alert handlers and predict the root cause
category, achieving up to 76.6\% accuracy on real incident
data~\cite{chen2024rcacopilot}. Ahmed et al.\ conduct the first large-scale study of GPT-3.x
models for recommending root causes and mitigation steps from incident reports at Microsoft,
finding that LLMs can surface relevant diagnostic context even without task-specific
fine-tuning~\cite{ahmed2023rcarecommend}. Roy et al.\ explore the use of LLM-based
ReAct~\cite{yao2023reactsynergizingreasoningacting} agents equipped with retrieval tools for RCA
on an out-of-distribution dataset of production incidents, showing that agentic tool use can
improve diagnostic coverage but also identifying key failure modes when retrieval returns
irrelevant context~\cite{roy2024llmagentRCA}. AetherLog, which is the most directly related
prior work to this thesis, integrates an LLM with a knowledge graph of historical incidents for
log-based RCA, demonstrating that structured retrieval of similar past cases substantially
reduces unsupported speculation and improves the precision of generated
diagnoses~\cite{cui2025aetherlog}.

Despite these promising results, single-LLM and single-agent RCA approaches share a set of
structural limitations that have not been fully resolved. The most significant is the tendency
toward hallucination: when the prompt does not contain sufficient or unambiguous evidence, an LLM
will produce a plausible-sounding but unsupported diagnosis, and there is no internal mechanism
to detect or correct this~\cite{du2023improvingfactualityreasoninglanguage}. This is particularly
consequential in RCA because an incorrect diagnosis can lead an engineer to perform unnecessary
or harmful remediation actions. A related limitation is the absence of hypothesis diversity: a
single model pass commits to one explanation without systematically considering alternatives. In
cases where evidence is ambiguous---which is common when failures propagate across multiple
components---premature commitment to a single hypothesis risks anchoring the operator on the
wrong causal narrative. Furthermore, single-pass generation provides no calibrated confidence
signal: the model's expressed confidence correlates poorly with actual correctness, making it
difficult for operators to decide when to trust the automated diagnosis and when to escalate
to manual investigation~\cite{chen2024rcacopilot,roy2024llmagentRCA}.

\section{Addressing the Gap: Agent Debate with Knowledge Graph Grounding}

This thesis proposes \systemname{}, a multi-agent, knowledge-graph-guided RCA system designed
to address the limitations of single-agent LLM approaches. The core insight is that the
reliability problems of single-agent RCA stem from two compounding factors: insufficient
evidence grounding and the absence of adversarial verification. We address both simultaneously.

Evidence grounding is improved by integrating a Neo4j-based knowledge graph (KG) that stores
historical incidents, their involved entities, and their diagnosed root causes. At diagnosis time,
the KG is queried for historical incidents that share entities with the current incident, and the
retrieved context is provided to the reasoning agents as structured external memory. This design
is motivated by the success of retrieval-augmented generation for knowledge-intensive
tasks~\cite{lewis2021retrievalaugmentedgenerationknowledgeintensivenlp} and by the specific
observation that operational incident history---when organized as a structured
graph~\cite{hogan2021knowledge}---provides precisely the kind of entity-centric, relation-aware
context that supports reliable causal inference.

Adversarial verification is implemented through a structured debate protocol in which three
specialized reasoning agents---a log-focused agent, a KG-focused agent, and a hybrid
agent---independently generate competing hypotheses. A dedicated judge agent then evaluates all
hypotheses against a fixed rubric covering evidence quality, reasoning strength, confidence
calibration, completeness, and consistency. This design is motivated by recent empirical evidence
that multi-agent debate substantially improves the factual accuracy of LLM-generated
content~\cite{du2023improvingfactualityreasoninglanguage}, and by the observation that
specialization and decomposition reduce the chance that all agents share the same blind
spots~\cite{wu2023autogen}. By forcing hypotheses to compete under an explicit scoring
mechanism, the system reduces the risk of majority-class collapse and overconfident but
unsupported diagnoses that are the primary failure modes of single-agent baselines.

\section{Research Questions}

The following research questions structure the experimental investigation in this thesis:

\begin{description}
  \item[RQ1:] Does a multi-agent debate approach achieve higher diagnostic accuracy and better
              class balance than single-agent LLM baselines for log-based RCA?
  \item[RQ2:] Does structured debate with explicit judge scoring reduce overconfident or
              unsupported predictions compared to single-pass generation?
  \item[RQ3:] Does knowledge graph integration, as an external grounding source, measurably
              improve diagnosis quality beyond what retrieval-augmented single-agent generation
              achieves?
  \item[RQ4:] Do the generated explanations provide sufficient evidence citation and causal
              reasoning to support operator triage?
\end{description}

\section{Contributions}

This thesis makes the following contributions:

\paragraph{Multi-agent RCA architecture.}
We design and implement \systemname{}, a six-agent RCA system comprising a Log Parser, a KG
Retrieval Agent, three specialized Reasoners (log-focused, KG-focused, and hybrid), and a Judge
Agent. Each agent has a well-defined role and operates over structured input/output contracts,
enabling systematic composition and debugging.

\paragraph{Structured debate protocol.}
We propose a multi-round debate protocol in which three reasoning agents independently generate
competing hypotheses, a judge evaluates them against a fixed five-criterion rubric, and agents
refine their hypotheses based on judge feedback in subsequent rounds. This protocol operationalizes
reliability as a measurable, evidence-based selection process rather than an informal heuristic.

\paragraph{Knowledge graph for operational memory.}
We implement a Neo4j-based incident knowledge graph that supports entity-centric retrieval of
similar historical incidents. The KG is populated incrementally during evaluation, simulating
an operational deployment where diagnostic history accumulates over time.

\paragraph{Empirical evaluation across three datasets.}
We evaluate \systemname{} on three publicly available log datasets spanning two domains:
Hadoop1 and HDFS\_v1 from LogHub~\cite{he2020loghub}, and the CMCC OpenStack failure
dataset~\cite{cui2025aetherlog}. Across all three datasets, the multi-agent system consistently
outperforms single-agent and RAG baselines in macro-averaged \fone{} score, with the most
substantial improvement on CMCC, where single-agent baselines collapse to uninformative
predictions.

\paragraph{Fully local and reproducible implementation.}
The entire system runs on local infrastructure using Ollama for LLM inference and Neo4j Community
Edition for graph storage, with no dependence on commercial API availability. All prompts,
intermediate outputs, and evaluation scripts are logged, enabling full reproduction of results.

\section{Thesis Organization}

Chapter~2 reviews the background and related work, covering traditional and ML-based RCA,
knowledge graphs for operational reasoning, LLM foundations from transformers to instruction
tuning and alignment, and multi-agent reasoning systems. Chapter~3 presents the system design
and architecture of \systemname{}, detailing the six agents, the knowledge graph schema, and the
debate protocol. Chapter~4 describes the experimental methodology, including label normalization,
baseline configurations, and evaluation metrics. Chapter~5 covers the implementation, including
local LLM inference, structured output handling, and the knowledge graph pipeline.
Chapter~6 presents the experimental results on all three datasets with detailed per-class
analysis and a misclassification audit. Chapter~7 discusses the findings, limitations, and threats
to validity. Chapter~8 concludes with a summary of contributions, answers to the research
questions, and directions for future work. Appendix~A provides representative case studies
extracted from system outputs.

% -*- coding: utf-8 -*-

\chapter{Background and Related Work}

\section{Log Parsing and Template Mining}

Raw log messages are free-form text that combine natural language with runtime variable tokens
such as IP addresses, block identifiers, and numeric values. Before any pattern-based or
learning-based analysis can be applied, these messages must be transformed into structured event
templates by masking the variable portions. This task is known as log parsing, and it serves as a
prerequisite to virtually all downstream log analytics
pipelines~\cite{zhu2019tools}.

Several parsing approaches have been proposed with different tradeoffs between accuracy, throughput,
and adaptability to new log formats. Drain organizes discovered templates in a fixed-depth parse
tree, enabling efficient online matching and incremental template
updates~\cite{he2017drain}. Its fixed-depth design ensures that the parse tree remains compact
even when the number of templates grows, making it well-suited for continuous log streams. Spell
takes a streaming approach, using longest common subsequence matching to maintain message types as
new logs arrive without requiring global re-clustering~\cite{du2016spell}. Earlier work such as
IPLoM employs iterative partitioning, clustering messages by progressively splitting on token
positions~\cite{makanju2009iplom}. A systematic benchmark of parsing tools confirms that different
parsers exhibit substantially different accuracy and efficiency profiles on different log
corpora~\cite{zhu2019tools}, making parser selection a practically important preprocessing decision.
In our system, the Log Parser Agent performs a semantic variant of this task using an LLM, which
allows it to handle heterogeneous log formats without requiring a pre-defined template vocabulary.

\section{Root Cause Analysis: Traditional Approaches}

It is important to distinguish log-based RCA from the closely related but distinct task of anomaly
detection. Anomaly detection determines whether a sequence of log events or metrics deviates from
a learned or rule-defined model of normal behavior; it produces a binary or scored label
characterizing whether an execution window is anomalous~\cite{du2017deeplog,zhang2019logrobust}.
RCA, by contrast, is concerned with explaining \emph{why} an incident occurred: which component or
sub-system is the root cause, how the fault propagated, and what remediation is appropriate.
RCA takes anomaly detection as a possible input signal---it may reason over sequences of anomalous
events---but its objective is causal attribution rather than deviation
scoring~\cite{meng2020localizing,wu2020microrca}. This thesis is concerned exclusively with the
RCA task; anomaly detection methods are reviewed here only to the extent that they inform the log
representation strategies used by our parser.

Traditional automated RCA approaches can be broadly categorized into two families. The first family
employs causal graph methods: a dependency graph among services or metrics is constructed, and
statistical or graph-theoretic algorithms identify the component whose anomalous behavior is most
causally associated with the observed incident. CauseInfer builds a hierarchical causality graph
from service performance metrics to support automatic fault
localization~\cite{chen2014causeinfer}. MicroRCA uses a service call graph weighted by anomalous
performance indicators to localize root causes in microservice
systems~\cite{wu2020microrca}. CloudRanger applies Bayesian network learning to system metric
data to identify the most likely causal node in the dependency
graph~\cite{wang2018cloudranger}. The localization accuracy of these methods depends critically
on the quality of the dependency model; when the model is incomplete or stale, fault propagation
paths may not be correctly captured. Moreover, constructing and maintaining accurate causal graphs
requires ongoing engineering effort and deep system expertise, which limits practical deployment
in large, rapidly evolving cloud environments.

The second family applies supervised machine learning to log-derived representations. Early work
by Xu et al.\ demonstrates that principal component analysis over log message count matrices can
detect and localize problems in large-scale HDFS deployments, establishing logs as a viable signal
for automated system problem detection~\cite{xu2009consolelogs}. Subsequent work trains sequence
models over event template sequences to recognize fault signatures: DeepLog uses LSTM networks to
model normal execution sequences and flags out-of-distribution events as
anomalous~\cite{du2017deeplog}. While these methods can achieve strong detection rates on
controlled benchmarks, they require labeled training data from the target system, and their
performance degrades when log statement formatting changes across software versions---a pervasive
practical problem demonstrated empirically by LogRobust~\cite{zhang2019logrobust}. Furthermore,
detection-oriented classifiers produce a fault label or anomaly score but do not generate an
interpretable causal explanation, which is precisely what incident response workflows
require~\cite{chen2024rcacopilot}.

\section{Datasets and Benchmarks}

Public benchmarks have been essential for enabling reproducible evaluation and comparison across
log analysis approaches. LogHub aggregates log data from more than fifteen distinct software
systems---including Hadoop, HDFS, OpenStack, Apache, and Android---and has become the standard
reference for both anomaly detection and RCA research~\cite{he2020loghub}. The Hadoop1 and
HDFS\_v1 datasets used in our evaluation are drawn from LogHub. The HDFS dataset traces back to
foundational work demonstrating that console log analysis can scale to detecting problems across
tens of thousands of nodes~\cite{xu2009consolelogs}. For cloud management platforms, the CMCC
dataset provides industrial OpenStack failure traces in which failures propagate across interacting
services such as Nova, Neutron, and RabbitMQ; we use this dataset to evaluate cross-domain
generalization to real enterprise incidents~\cite{cui2025aetherlog}.

\section{Knowledge Graphs for Operational Reasoning}

Knowledge graphs represent entities, their attributes, and the relations between them in a
structured, queryable form. As a general technology, they support entity-centric retrieval,
multi-hop relational reasoning, and the integration of heterogeneous evidence
sources~\cite{hogan2021knowledge}. These properties are directly useful in AIOps: a KG that
stores historical incidents, the entities they involved, and their diagnosed root causes provides
a structured operational memory that can inform diagnosis of new incidents whose entities overlap
with past cases.

AetherLog is the most directly related prior work that combines KG retrieval with LLM-based
reasoning for log-based RCA~\cite{cui2025aetherlog}. By querying a KG of historical incidents and
injecting the retrieved context into the LLM prompt, AetherLog demonstrates that structured
historical grounding substantially reduces unsupported speculation and improves precision compared
to a purely parametric LLM baseline. Our work extends this direction by embedding KG retrieval
within a multi-agent debate framework, so that the retrieved context is evaluated critically
rather than accepted uncritically by a single reasoning pass.

\section{Retrieval-Augmented Generation}

Retrieval-augmented generation (RAG) augments a language model's parametric knowledge with
evidence retrieved from an external corpus at inference
time~\cite{lewis2021retrievalaugmentedgenerationknowledgeintensivenlp}. For knowledge-intensive
tasks where the relevant information may not have been adequately captured during pretraining, RAG
substantially improves factual accuracy by conditioning generation on retrieved documents. In the
RCA setting, retrieved evidence takes the form of incident reports, historical log summaries, or
structured KG subgraphs. The RAG paradigm motivates our KG Retrieval Agent, which supplies the
reasoning agents with entity-matched historical context. However, a limitation of standard RAG
pipelines is that retrieval quality is not verified: when retrieved context is irrelevant or
misleading, a single-agent model can still produce a confident but incorrect diagnosis. Roy et
al.\ explicitly identify this as a failure mode of LLM-based RCA agents, finding that
out-of-distribution retrieval can actively harm diagnostic
accuracy~\cite{roy2024llmagentRCA}. Our debate protocol addresses this by requiring the judge to
score evidence quality, penalizing hypotheses that cite weak or irrelevant retrieved facts.

\section{LLM-Based RCA: Current State of the Art}

The application of large language models to automated RCA has emerged as an active research
direction over the past two years, motivated by the observation that LLMs can interpret
heterogeneous log text and generate natural-language explanations without requiring
dataset-specific training or hand-crafted
rules~\cite{brown2020languagemodelsfewshotlearners,ouyang2022traininglanguagemodelsfollow}.

RCACopilot, deployed at Microsoft on Azure cloud incidents, represents the most prominent
production-scale LLM-based RCA system to date~\cite{chen2024rcacopilot}. It employs per-alert
incident handlers that collect diagnostic data from logs, traces, and metrics, and feeds this
aggregated context to an LLM for root cause category prediction. Evaluated on a year's worth of
production incidents, RCACopilot achieves up to 76.6\% RCA accuracy and demonstrates that LLMs
can match or exceed the accuracy of rule-based on-call systems when given sufficiently curated
diagnostic inputs. A key design choice is the separation of data aggregation (handled by
specialized handlers per alert type) from reasoning (handled by the LLM), which reduces the
context burden on the model. However, RCACopilot operates as a single-agent system with no
cross-validation of its prediction, and its reliance on specialized per-alert handlers limits
generalizability to new alert types without additional engineering.

Ahmed et al.\ conduct an early large-scale empirical study of GPT-3.x models for recommending
root causes and mitigation steps from Microsoft incident reports, finding that LLM outputs are
useful even without fine-tuning but that their quality degrades substantially for incidents whose
diagnostic signals do not appear in the model's training
distribution~\cite{ahmed2023rcarecommend}. This finding motivates the use of retrieval-based
grounding to supply in-context evidence that the model cannot rely on memorized knowledge to
provide.

Roy et al.\ study LLM-based ReAct agents equipped with retrieval tools for RCA on an
out-of-distribution dataset of production incidents~\cite{roy2024llmagentRCA}. Their evaluation
reveals a fundamental tension: agentic tool use can improve diagnostic coverage by retrieving
relevant evidence, but it also introduces new failure modes when the retrieval system returns
irrelevant results that the agent uncritically incorporates into its reasoning. This finding
directly informs our design: rather than relying on a single agent to both retrieve and reason,
we separate the retrieval function into a dedicated KG Retrieval Agent and subject the reasoning
outputs to competitive evaluation by the judge.

AetherLog specifically targets log-based RCA and integrates KG retrieval within the
LLM-based diagnostic pipeline~\cite{cui2025aetherlog}. Unlike the cloud incident systems of
RCACopilot and Ahmed et al., which consume pre-processed diagnostic bundles, AetherLog operates
on raw log sequences, making it the most architecturally comparable prior work to ours. Our
system extends AetherLog's core insight---that structured historical grounding improves LLM-based
log diagnosis---by adding the multi-agent debate and judging layers that AetherLog lacks.

\section{Foundations of Large Language Models}

Understanding why LLMs are effective for log-based reasoning requires reviewing the architectural
and training advances that underpin them.

\subsection{Transformer architectures}

The transformer architecture replaces recurrence with token-level self-attention, enabling parallel
computation over sequences and flexible modeling of long-range
dependencies~\cite{vaswani2023attentionneed}. This architectural property is directly relevant
to log analysis: failure signatures are often sparse, with the critical error token appearing many
lines after the triggering event, and self-attention provides a mechanism to associate these
distant tokens without the vanishing-gradient problems of recurrent models. BERT demonstrates that
deep bidirectional transformers pretrained on large corpora learn powerful contextual
representations applicable to a wide range of downstream
tasks~\cite{devlin2019bertpretrainingdeepbidirectional}, while RoBERTa establishes best practices
for pretraining data and hyperparameter choices that substantially improve representation
quality~\cite{liu2019robertarobustlyoptimizedbert}. Autoregressive decoder models such as
GPT-3 demonstrate that sufficiently large transformers can perform complex reasoning and
multi-step generation in a few-shot prompting setting, without any task-specific gradient
updates~\cite{brown2020languagemodelsfewshotlearners}. For RCA, this means that a capable
decoder LLM can generate structured hypotheses, evidence lists, and remediation narratives in
a single forward pass, conditioned solely on the prompt content.

\subsection{Scaling laws and emergent capabilities}

Empirical scaling laws establish that language model performance improves predictably as a
power-law function of model size, dataset size, and training
compute~\cite{kaplan2020scalinglawsneurallanguage}. Chinchilla scaling analysis further refines
this understanding, arguing that most existing LLMs were under-trained relative to their size and
that compute-optimal training requires substantially more data per parameter than earlier
recipes~\cite{hoffmann2022trainingcomputeoptimallargelanguage}. Large-scale training experiments
such as Gopher and PaLM characterize how emergent capabilities---including multi-step reasoning,
code generation, and in-context learning---appear at scale and improve non-monotonically with
model size~\cite{rae2022scalinglanguagemodelsmethods,chowdhery2022palmscalinglanguagemodeling}.
These findings have practical implications for our work: they explain why modern 7B-class open
models, which benefit from compute-efficient training recipes derived from scaling research, can
exhibit strong semantic generalization over heterogeneous operational text even in a zero-shot
setting.

\subsection{Instruction tuning and alignment}

Raw pretrained LLMs generate fluent text but do not reliably follow task instructions, adhere to
output format constraints, or avoid overconfident assertions. Instruction tuning and alignment
methods are designed to instill these behaviors. InstructGPT demonstrates that training with
reinforcement learning from human feedback (RLHF)---using PPO to optimize against a learned
reward model of human preferences---substantially improves instruction-following quality compared
to supervised fine-tuning
alone~\cite{ouyang2022traininglanguagemodelsfollow,schulman2017proximalpolicyoptimizationalgorithms}.
Complementary instruction-tuning approaches demonstrate that scale is not the only route to
zero-shot task generalization: finetuned language models evaluated on held-out task types show
strong zero-shot performance when instruction coverage in fine-tuning is
broad~\cite{wei2022finetunedlanguagemodelszeroshot}. Self-Instruct shows that self-generated
instructions can bootstrap alignment data at low cost~\cite{wang2023selfinstructaligninglanguagemodels},
while LIMA demonstrates that a small, carefully curated set of alignment examples can be
surprisingly effective~\cite{zhou2023limaalignment}. Direct preference optimization (DPO) offers
a simplified alternative to RLHF that avoids explicit reward model training by directly optimizing
the language model on preference
data~\cite{rafailov2024directpreferenceoptimizationlanguage}.

In the context of RCA, alignment is practically important in three respects. First, the system
must produce outputs in a strict JSON schema; misformatted outputs cannot be parsed and disrupt
the pipeline at scale. Second, the judge agent must produce calibrated scores rather than simply
asserting that one hypothesis is ``better.'' Third, reasoning agents must acknowledge uncertainty
rather than generating overconfident diagnoses when evidence is weak. Modern instruction-tuned
LLMs provide a strong foundation for these requirements, though reliability for operational tasks
still requires the additional architectural guardrails that our debate protocol provides.

\subsection{Parameter-efficient adaptation}

Even well-aligned foundation models may not have been exposed to the specific log terminology,
error patterns, and operational conventions of a target system during pretraining. Parameter-efficient
fine-tuning (PEFT) provides a practical mechanism for domain
adaptation~\cite{raffel2023exploringlimitstransferlearning}: LoRA inserts low-rank decomposition
matrices into transformer layers, adapting model behavior with a fraction of the parameters and
compute required for full fine-tuning~\cite{hu2021loralowrankadaptationlarge}. QLoRA further
reduces the resource footprint by combining quantization of the base model with LoRA adapters,
enabling fine-tuning of 7B-class models on commodity
hardware~\cite{dettmers2023qloraefficientfinetuningquantized}. While we do not perform
domain-specific fine-tuning in this thesis, these methods represent the most practical path for
future operational deployments where privacy constraints require on-premise adaptation on
organization-specific incident corpora.

\section{LLM Reasoning Strategies and Multi-Agent Verification}

Even a well-aligned, capable LLM can produce incorrect diagnoses when evidence is ambiguous or
incomplete. A growing body of work addresses this reliability problem through structured reasoning
strategies and multi-agent architectures.

Chain-of-thought prompting demonstrates that instructing a model to generate intermediate reasoning
steps before producing a final answer substantially improves performance on multi-step reasoning
tasks~\cite{wei2023chainofthoughtpromptingelicitsreasoning}. The key mechanism is that explicit
intermediate steps anchor the model's attention on the relevant sub-problems and reduce the chance
of shortcut reasoning that bypasses the evidentiary chain. Self-consistency extends this idea by
sampling multiple reasoning traces from the same model and selecting the most frequent final answer,
thereby reducing sensitivity to any single sampled
trajectory~\cite{wang2023selfconsistencyimproveschainthought}. Large LLMs are also zero-shot
reasoners~\cite{kojima2023largelanguagemodelszeroshot}, meaning that step-by-step reasoning can
be elicited simply by appending ``Let's think step by step'' to a prompt. ReAct interleaves
reasoning with tool invocation, enabling models to alternate between deliberation and action in
interactive environments~\cite{yao2023reactsynergizingreasoningacting}; Roy et al.\ apply this
approach to RCA specifically~\cite{roy2024llmagentRCA}.

Multi-agent debate takes a different approach to reliability: rather than improving the reasoning
of a single model, it introduces multiple agents that independently generate answers and then
critique each other's responses. Du et al.\ provide the first systematic empirical study of
multi-agent debate, demonstrating that having multiple LLM instances propose and refine answers
through iterative debate reduces factual errors and improves reasoning
quality~\cite{du2023improvingfactualityreasoninglanguage}. AutoGen formalizes multi-agent
conversation and tool use within a general framework that supports heterogeneous agent roles,
custom conversation topologies, and human-in-the-loop
interaction~\cite{wu2023autogen}. These frameworks establish the architectural vocabulary that
we draw on in designing our debate protocol, adapted to the specific requirements of log-based
RCA: hypothesis generation must be grounded in log evidence and historical KG context, the judge
must score against an operational rubric, and the protocol must terminate within a bounded number
of rounds to remain practical.

\section{Summary and Research Gap}

The literature reviewed in this chapter reveals a clear progression from template-based log
parsers through statistical RCA methods to LLM-based diagnostic systems, and from single-agent
LLM approaches to agent frameworks with tool use and retrieval augmentation. At the same time,
it reveals a gap that this thesis addresses.

Traditional causal-graph and supervised ML approaches to RCA require either a well-maintained
dependency model or labeled training data from the target system, neither of which is reliably
available in cross-domain or rapidly evolving cloud
environments~\cite{zhang2024cloudRCAsurvey,zhang2019logrobust}. LLM-based RCA systems such as
RCACopilot and AetherLog demonstrate that LLMs can diagnose incidents without task-specific
training, but they share the structural limitation of single-agent, single-pass generation: there
is no mechanism to verify whether the produced hypothesis is supported by the available evidence
or to systematically explore alternative explanations before committing to a
diagnosis~\cite{chen2024rcacopilot,roy2024llmagentRCA}. Knowledge graph grounding improves
factuality by supplying external evidence, but retrieval alone does not prevent the model from
ignoring or misinterpreting the retrieved context~\cite{cui2025aetherlog}.

Multi-agent debate has been shown to improve factual accuracy and reasoning quality in general
LLM settings~\cite{du2023improvingfactualityreasoninglanguage}, but this approach has not
previously been combined with knowledge graph grounding and applied to log-based RCA. This thesis
fills that gap by proposing a system in which multi-agent hypothesis generation, KG-based
evidence retrieval, and explicit judge-based evaluation are integrated into a single coherent
pipeline, and by evaluating this design empirically on three publicly available incident
datasets.

\include{chapters/03_system_design}
\include{chapters/03b_methodology}
\include{chapters/04_implementation}
% -*- coding: utf-8 -*-

\chapter{Experimental Setup and Results}

\section{Datasets}
We evaluate on three datasets:
\begin{itemize}
  \item \textbf{Hadoop1 (LogHub)}: 55 labeled applications with four fault types (normal,
        machine\_down, network\_disconnection, disk\_full).
  \item \textbf{CMCC (LogKG)}: 93 OpenStack failure cases with 7 failure types.
  \item \textbf{HDFS\_v1 (LogHub)}: 575,061 block traces with binary labels; we evaluate a
        balanced sample of 200 blocks.
\end{itemize}

\section{Evaluation Schemes and Metrics}
For Hadoop1, we evaluate using both \textbf{strict} (4-class) and \textbf{coarse} (3-class) label
schemes. We report accuracy and macro-\fone{}; for imbalanced settings, macro-\fone{} is
emphasized.

\subsection{Metrics}
We report the following metrics:
\begin{itemize}
  \item \textbf{Accuracy}: the fraction of cases whose predicted label matches the ground truth.
  \item \textbf{Precision/Recall}: per-class precision and recall, aggregated as \textbf{macro
        averages} (unweighted mean across classes).
  \item \textbf{Macro-\fone{}}: the unweighted mean of per-class \fone{} scores.
\end{itemize}
Macro-averaged metrics are important in operational settings because they reduce the dominance of
frequent classes and better reflect performance on rare but operationally critical failures.

\subsection{Strict vs.\ coarse label schemes (Hadoop1)}
Hadoop1 supports two evaluation views. The strict scheme uses the four original labels, while the
coarse scheme merges semantically similar connectivity-related faults into a single class. This
provides a more operator-aligned view because multiple low-level causes (e.g., machine down vs.\
network disconnection) can manifest as connectivity disruption at the application level.

\subsection{LLM-based evaluation protocol}
All three pipelines produce natural-language hypotheses and must be mapped to a dataset label for
evaluation. This introduces two practical constraints. First, LLM outputs must be sufficiently
structured to enable automatic parsing at scale. Second, the evaluation must avoid
over-interpreting ambiguous text. For this reason, we prompt models for structured JSON outputs
and perform conservative label normalization.

This setup follows a growing trend in using transformer-based LLMs for complex reasoning and
explanation generation~\cite{vaswani2023attentionneed,brown2020languagemodelsfewshotlearners}.
However, reliability is not guaranteed by the foundation model alone, and instruction-following
behavior can vary with model family and alignment
training~\cite{ouyang2022traininglanguagemodelsfollow}. Our multi-agent protocol addresses this by
generating diverse candidates and applying an explicit judge for evidence-based
selection~\cite{du2023improvingfactualityreasoninglanguage,wu2023autogen}.

\subsection{Baselines}
We compare against:
\begin{itemize}
  \item \textbf{Single-agent}: a single LLM call that reasons over incident logs.
  \item \textbf{RAG}: a single LLM call augmented with retrieved historical context. RAG is a
        common mechanism for improving factuality by grounding generation in external
        evidence~\cite{lewis2021retrievalaugmentedgenerationknowledgeintensivenlp}.
  \item \textbf{Logistic Regression (LR)}: a supervised ML baseline using Drain-parsed log
        template event-count feature vectors with scikit-learn logistic regression, following
        the Loglizer evaluation
        protocol~\cite{he2016loglizer}. We use a stratified 70/30 train-test split for
        Hadoop1 and CMCC; for HDFS we report numbers from~\cite{he2016loglizer} directly.
        Unlike the LLM pipelines, LR requires labeled training data and cannot generalize
        zero-shot across domains.
\end{itemize}

\section{Hadoop1 Experimental Setup}

\subsection{Dataset Statistics}

\begin{table}[H]
\centering
\caption{Hadoop1 Dataset Statistics}
\label{tab:hadoop1-dataset}
\begin{tabular}{lrrl}
\toprule
\textbf{Label} & \textbf{Count} & \textbf{Percentage} & \textbf{Description} \\
\midrule
normal                 & 11 & 20.0\%  & No injected fault \\
machine\_down          & 28 & 50.9\%  & Node failure injected \\
network\_disconnection & 7  & 12.7\%  & Network partition injected \\
disk\_full             & 9  & 16.4\%  & Disk space exhaustion \\
\midrule
\textbf{Total}         & \textbf{55} & \textbf{100\%} & \\
\bottomrule
\end{tabular}
\end{table}

\begin{table}[H]
\centering
\caption{Hadoop1 Log Volume Statistics}
\label{tab:hadoop1-volume}
\begin{tabular}{lr}
\toprule
\textbf{Metric} & \textbf{Value} \\
\midrule
Total Log Entries           & 394,308 \\
Total Log Files             & 978     \\
Applications                & 55      \\
Avg.\ Lines per Application & 7,169   \\
Dataset Size                & 49\,MB  \\
\bottomrule
\end{tabular}
\end{table}

\subsection{System Comparison}

\begin{table}[H]
\centering
\caption{System Comparison}
\label{tab:systems}
\begin{tabular}{llll}
\toprule
\textbf{Aspect}         & \textbf{Multi-Agent}    & \textbf{RAG}          & \textbf{Single-Agent} \\
\midrule
LLM Calls per app       & 8--15                   & 1                     & 1 \\
Perspectives            & 3 (Log, KG, Hybrid)     & 1                     & 1 \\
KG Integration          & Yes                     & Yes (retrieval only)  & No \\
Debate/Refinement       & Yes (2--3 rounds)       & No                    & No \\
Cross-Validation        & Yes (Judge)             & No                    & No \\
Runtime per app         & 3--5 minutes            & 30--45 seconds        & 30 seconds \\
\bottomrule
\end{tabular}
\end{table}

\section{Hadoop1 Results}

\subsection{Multi-Agent Pipeline Results}

\subsubsection{Strict (4-class) Evaluation}

\begin{table}[H]
\centering
\caption{Multi-Agent Strict Evaluation Results}
\label{tab:ma-strict}
\begin{tabular}{lr}
\toprule
\textbf{Metric} & \textbf{Value} \\
\midrule
Accuracy          & 21.8\% (12/55) \\
Macro Precision   & 41.7\% \\
Macro Recall      & 30.6\% \\
Macro \fone{}     & 21.6\% \\
Unknown Predictions & 9 \\
\bottomrule
\end{tabular}
\end{table}

\begin{table}[H]
\centering
\caption{Multi-Agent Per-Class Results (Strict)}
\label{tab:ma-strict-perclass}
\begin{tabular}{lrrrr}
\toprule
\textbf{Class} & \textbf{Support} & \textbf{Precision} & \textbf{Recall} & \textbf{\fone{}} \\
\midrule
normal                 & 11 & 0.0\%   & 0.0\%   & 0.0\%  \\
machine\_down          & 28 & 50.0\%  & 14.3\%  & 22.2\% \\
network\_disconnection & 7  & 16.7\%  & 85.7\%  & 27.9\% \\
disk\_full             & 9  & 100.0\% & 22.2\%  & 36.4\% \\
\bottomrule
\end{tabular}
\end{table}

\begin{table}[H]
\centering
\caption{Multi-Agent Confusion Matrix (Strict)}
\label{tab:ma-cm-strict}
\begin{tabular}{l|ccccc}
\toprule
& \multicolumn{5}{c}{\textbf{Predicted}} \\
\textbf{Ground Truth} & normal & mach\_down & net\_disc & disk\_full & unknown \\
\midrule
normal (n=11)       & 0 & 2 & 6 & 0 & 3 \\
machine\_down (n=28) & 0 & 4 & 21 & 0 & 3 \\
network\_disc (n=7)  & 0 & 1 & 6 & 0 & 0 \\
disk\_full (n=9)     & 0 & 1 & 3 & 2 & 3 \\
\bottomrule
\end{tabular}
\end{table}

\paragraph{Interpretation (strict).}
The strict scheme exposes a mismatch between symptom-level evidence and injected-fault labels.
The confusion matrix shows that many \texttt{machine\_down} and \texttt{normal} cases are predicted
as \texttt{network\_disconnection}, suggesting that observable log cues in Hadoop1 frequently
emphasize connectivity symptoms. The multi-agent pipeline produces a non-trivial number of
\texttt{unknown} predictions, which occurs when the free-text hypothesis does not map cleanly to
one of the strict labels.

The per-class table highlights a common operational phenomenon: a model can have strong precision
on a rare class (e.g., \texttt{disk\_full}) when it emits the label, but low recall if it rarely
selects it. This motivates reporting macro metrics alongside accuracy.

\subsubsection{Coarse (3-class) Evaluation}

\begin{table}[H]
\centering
\caption{Multi-Agent Coarse Evaluation Results}
\label{tab:ma-coarse}
\begin{tabular}{lr}
\toprule
\textbf{Metric} & \textbf{Value} \\
\midrule
Accuracy            & 61.8\% (34/55) \\
Macro Precision     & 57.6\% \\
Macro Recall        & 37.9\% \\
Macro \fone{}       & 39.1\% \\
Unknown Predictions & 9 \\
\bottomrule
\end{tabular}
\end{table}

\begin{table}[H]
\centering
\caption{Multi-Agent Per-Class Results (Coarse)}
\label{tab:ma-coarse-perclass}
\begin{tabular}{lrrrr}
\toprule
\textbf{Class} & \textbf{Support} & \textbf{Precision} & \textbf{Recall} & \textbf{\fone{}} \\
\midrule
normal       & 11 & 0.0\%   & 0.0\%   & 0.0\%  \\
connectivity & 35 & 72.7\%  & 91.4\%  & 81.0\% \\
disk\_full   & 9  & 100.0\% & 22.2\%  & 36.4\% \\
\bottomrule
\end{tabular}
\end{table}

\begin{table}[H]
\centering
\caption{Multi-Agent Confusion Matrix (Coarse)}
\label{tab:ma-cm-coarse}
\begin{tabular}{l|cccc}
\toprule
& \multicolumn{4}{c}{\textbf{Predicted}} \\
\textbf{Ground Truth} & normal & connectivity & disk\_full & unknown \\
\midrule
normal (n=11)       & 0 & 8  & 0 & 3 \\
connectivity (n=35) & 0 & 32 & 0 & 3 \\
disk\_full (n=9)    & 0 & 4  & 2 & 3 \\
\bottomrule
\end{tabular}
\end{table}

\paragraph{Interpretation (coarse).}
The coarse scheme reduces label ambiguity by aligning evaluation with the dominant symptom-level
evidence present in Hadoop1 logs. As a result, connectivity-related cases become easier to
classify and macro-\fone{} improves compared to the strict scheme. This is consistent with prior
findings that label semantics and preprocessing choices can strongly affect reported
performance~\cite{le2022logdlhowfar}.

\subsection{Single-Agent Baseline Results}

\subsubsection{Strict (4-class) Evaluation}

\begin{table}[H]
\centering
\caption{Single-Agent Strict Evaluation Results}
\label{tab:sa-strict}
\begin{tabular}{lr}
\toprule
\textbf{Metric} & \textbf{Value} \\
\midrule
Accuracy            & 50.9\% (28/55) \\
Macro Precision     & 13.2\% \\
Macro Recall        & 25.0\% \\
Macro \fone{}       & 17.3\% \\
Unknown Predictions & 0 \\
\bottomrule
\end{tabular}
\end{table}

\begin{table}[H]
\centering
\caption{Single-Agent Per-Class Results (Strict)}
\label{tab:sa-strict-perclass}
\begin{tabular}{lrrrr}
\toprule
\textbf{Class} & \textbf{Support} & \textbf{Precision} & \textbf{Recall} & \textbf{\fone{}} \\
\midrule
normal                 & 11 & 0.0\%  & 0.0\%   & 0.0\%  \\
machine\_down          & 28 & 52.8\% & 100.0\% & 69.1\% \\
network\_disconnection & 7  & 0.0\%  & 0.0\%   & 0.0\%  \\
disk\_full             & 9  & 0.0\%  & 0.0\%   & 0.0\%  \\
\bottomrule
\end{tabular}
\end{table}

\subsubsection{Coarse (3-class) Evaluation}

\begin{table}[H]
\centering
\caption{Single-Agent Coarse Evaluation Results}
\label{tab:sa-coarse}
\begin{tabular}{lr}
\toprule
\textbf{Metric} & \textbf{Value} \\
\midrule
Accuracy            & 61.8\% (34/55) \\
Macro Precision     & 21.4\% \\
Macro Recall        & 32.4\% \\
Macro \fone{}       & 25.8\% \\
Unknown Predictions & 0 \\
\bottomrule
\end{tabular}
\end{table}

\subsection{RAG Baseline Results}

\subsubsection{Strict (4-class) Evaluation}

\begin{table}[H]
\centering
\caption{RAG Baseline Strict Evaluation Results}
\label{tab:rag-strict}
\begin{tabular}{lr}
\toprule
\textbf{Metric} & \textbf{Value} \\
\midrule
Accuracy            & 49.1\% (27/55) \\
Macro Precision     & 29.4\% \\
Macro Recall        & 27.9\% \\
Macro \fone{}       & 24.6\% \\
Unknown Predictions & 0 \\
\bottomrule
\end{tabular}
\end{table}

\subsubsection{Coarse (3-class) Evaluation}

\begin{table}[H]
\centering
\caption{RAG Baseline Coarse Evaluation Results}
\label{tab:rag-coarse}
\begin{tabular}{lr}
\toprule
\textbf{Metric} & \textbf{Value} \\
\midrule
Accuracy            & 67.3\% (37/55) \\
Macro Precision     & 44.7\% \\
Macro Recall        & 40.7\% \\
Macro \fone{}       & 37.9\% \\
Unknown Predictions & 0 \\
\bottomrule
\end{tabular}
\end{table}

\subsection{Logistic Regression (LR) Baseline Results}

\subsubsection{Strict (4-class) Evaluation}

\begin{table}[H]
\centering
\caption{LR Baseline Strict Evaluation Results (supervised, 70/30 split)}
\label{tab:lr-strict}
\begin{tabular}{lr}
\toprule
\textbf{Metric} & \textbf{Value} \\
\midrule
Accuracy            & 52.7\% (29/55) \\
Macro Precision     & 36.1\% \\
Macro Recall        & 27.8\% \\
Macro \fone{}       & 26.1\% \\
\bottomrule
\end{tabular}
\end{table}

\subsubsection{Coarse (3-class) Evaluation}

\begin{table}[H]
\centering
\caption{LR Baseline Coarse Evaluation Results (supervised, 70/30 split)}
\label{tab:lr-coarse}
\begin{tabular}{lr}
\toprule
\textbf{Metric} & \textbf{Value} \\
\midrule
Accuracy            & 70.9\% (39/55) \\
Macro Precision     & 52.3\% \\
Macro Recall        & 42.6\% \\
Macro \fone{}       & 44.7\% \\
\bottomrule
\end{tabular}
\end{table}

\paragraph{Interpretation.}
Even with labeled training data, LR achieves only 26.1\% macro-\fone{} under the strict
4-class scheme. The class imbalance (machine\_down accounts for 50.9\% of cases) pushes LR
towards the majority class, mirroring the behaviour seen in single-agent LLM predictions.
Under the coarse scheme, LR improves to 44.7\% macro-\fone{}, outperforming the zero-shot
multi-agent system (39.1\%) on this dataset. This is expected: LR is a supervised method
trained on the same distribution, giving it an inherent advantage over zero-shot inference
that will not hold across domains.

\subsection{Comparative Summary and Key Findings}

\begin{table}[H]
\centering
\caption{Comparative Summary: All Approaches on Hadoop1.
  \textsuperscript{\dag}Supervised LR with 70/30 stratified split; not directly comparable to zero-shot LLM pipelines.}
\label{tab:comparison}
\begin{tabular}{lrrrr}
\toprule
\textbf{Metric} & \textbf{Multi-Agent} & \textbf{RAG} & \textbf{Single-Agent} & \textbf{LR\textsuperscript{\dag}} \\
\midrule
\multicolumn{5}{l}{\textit{Strict Evaluation (4-class)}} \\
Accuracy         & 21.8\% & 49.1\% & 50.9\% & 52.7\% \\
Macro Precision  & 41.7\% & 29.4\% & 13.2\% & 36.1\% \\
Macro Recall     & 30.6\% & 27.9\% & 25.0\% & 27.8\% \\
Macro \fone{}    & 21.6\% & 24.6\% & 17.3\% & 26.1\% \\
\midrule
\multicolumn{5}{l}{\textit{Coarse Evaluation (3-class)}} \\
Accuracy         & 61.8\% & 67.3\% & 61.8\% & 70.9\% \\
Macro Precision  & 57.6\% & 44.7\% & 21.4\% & 52.3\% \\
Macro Recall     & 37.9\% & 40.7\% & 32.4\% & 42.6\% \\
Macro \fone{}    & 39.1\% & 37.9\% & 25.8\% & \textbf{44.7\%} \\
\midrule
\multicolumn{5}{l}{\textit{Per-Class \fone{} (Coarse)}} \\
Normal       & 0.0\%  & 0.0\%  & 0.0\%  & 18.2\% \\
Connectivity & 81.0\% & 80.5\% & 77.3\% & \textbf{82.1\%} \\
Disk Full    & 36.4\% & 33.3\% & 0.0\%  & 33.9\% \\
\bottomrule
\end{tabular}
\end{table}

\paragraph{Strict vs.\ coarse takeaway.}
Across pipelines, strict Hadoop1 accuracy can be misleading: a model that over-predicts the
majority label can appear strong on accuracy while failing minority classes (macro-\fone{} near
zero for some classes). The multi-agent approach trades off strict accuracy for improved class
balance under the coarse scheme, which is often closer to diagnostic utility. The supervised
LR baseline achieves the highest coarse macro-\fone{} (44.7\%), but requires labeled in-domain
training data, making it unsuitable for zero-shot deployment across new environments.

\section{Cross-Dataset Validation}

\subsection{CMCC Dataset: OpenStack Multi-Class RCA}
\label{sec:cmcc-results}

\begin{table}[H]
\centering
\caption{CMCC Dataset Characteristics}
\label{tab:cmcc-dataset}
\begin{tabular}{lr}
\toprule
\textbf{Attribute} & \textbf{Value} \\
\midrule
Total Cases      & 93 \\
Domain           & OpenStack Cloud \\
Failure Classes  & 7 \\
Log Format       & Pre-parsed CSV \\
\bottomrule
\end{tabular}
\end{table}

\begin{table}[H]
\centering
\caption{CMCC Failure Type Distribution}
\label{tab:cmcc-distribution}
\begin{tabular}{lrl}
\toprule
\textbf{Failure Type} & \textbf{Count} & \textbf{Description} \\
\midrule
AMQP                        & 25 & RabbitMQ connection failures \\
Mysql                       & 18 & Database connection errors \\
CreateErrorFlavor           & 16 & Nova flavor creation errors \\
CreateErrorNovaConductor    & 13 & Nova conductor service errors \\
Down                        & 12 & Service unavailable \\
CreateErrorLinuxbridgeAgent &  9 & Neutron agent errors \\
\bottomrule
\end{tabular}
\end{table}

\begin{table}[H]
\centering
\caption{CMCC 7-Class RCA Results.
  \textsuperscript{\dag}Supervised LR with 70/30 stratified split; requires labeled training data.}
\label{tab:cmcc-results}
\begin{tabular}{lrrrr}
\toprule
\textbf{Pipeline} & \textbf{Accuracy} & \textbf{Macro \fone{}} & \textbf{Unknown} & \textbf{vs.\ Multi-Agent} \\
\midrule
\textbf{Multi-Agent} & \textbf{61.3\%} & \textbf{55.6\%} & 17.2\% & --- \\
RAG                  & 8.6\%           & 10.3\%          & 70.0\% & $-$45.3\,pp \\
Single-Agent         & 4.3\%           & 5.3\%           & 60.2\% & $-$50.3\,pp \\
LR\textsuperscript{\dag} & 40.9\%     & 33.5\%          & ---    & $-$22.1\,pp \\
\bottomrule
\end{tabular}
\end{table}

\paragraph{Interpretation.}
CMCC highlights the difficulty of mapping heterogeneous OpenStack logs into a discrete multi-class
taxonomy. The single-agent and RAG baselines show a very large fraction of \texttt{unknown}
predictions, indicating that a single-pass explanation often fails to align with the label space.
The multi-agent system reduces \texttt{unknown} outcomes by enforcing structured evaluation and
refinement: multiple hypotheses compete and the judge favors hypotheses with concrete evidence
and consistent causal narratives.

More broadly, transformer LMs can generalize semantically in a few-shot
setting~\cite{brown2020languagemodelsfewshotlearners}, but mapping free-text reasoning into closed
categories benefits from explicit verification and
judging~\cite{du2023improvingfactualityreasoninglanguage}.

\begin{table}[H]
\centering
\caption{CMCC Per-Class Performance (Multi-Agent)}
\label{tab:cmcc-perclass}
\begin{tabular}{lrrrr}
\toprule
\textbf{Class} & \textbf{Support} & \textbf{Precision} & \textbf{Recall} & \textbf{\fone{}} \\
\midrule
Mysql                       & 18 & 94.7\%  & 100.0\% & 97.3\% \\
AMQP                        & 25 & 71.0\%  & 88.0\%  & 78.6\% \\
CreateErrorNovaConductor    & 13 & 55.6\%  & 76.9\%  & 64.5\% \\
CreateErrorLinuxbridgeAgent &  9 & 100.0\% & 44.4\%  & 61.5\% \\
CreateErrorFlavor           & 16 & 100.0\% & 18.8\%  & 31.6\% \\
Down                        & 12 & 0.0\%   & 0.0\%   & 0.0\%  \\
\midrule
\textbf{Macro Average}      & 93 & 70.2\%  & 54.7\%  & 55.6\% \\
\bottomrule
\end{tabular}
\end{table}

\paragraph{Per-class behavior.}
Classes with distinctive error signatures (e.g., \texttt{Mysql}, \texttt{AMQP}) tend to be
predicted more reliably than classes that share symptom-level cues. This is consistent with the
intuition that log-only reasoning is strongest when error messages are explicit, and weaker when
failures manifest indirectly through downstream components.

\subsection{HDFS\_v1 Dataset: Large-Scale Binary Anomaly Detection}

\begin{table}[H]
\centering
\caption{HDFS\_v1 Dataset Characteristics}
\label{tab:hdfs-dataset}
\begin{tabular}{lr}
\toprule
\textbf{Attribute} & \textbf{Value} \\
\midrule
Total Block Traces & 575,061 \\
Normal Blocks      & 558,223 (97.1\%) \\
Anomaly Blocks     & 16,838  (2.9\%)  \\
Evaluation Sample  & 200 (balanced)   \\
Domain             & Hadoop HDFS      \\
Label Type         & Binary (Normal/Anomaly) \\
\bottomrule
\end{tabular}
\end{table}

\paragraph{Why HDFS\_v1 matters.}
HDFS-style benchmarks are closely related to early work showing that console logs contain enough
signal to detect system problems when mined
appropriately~\cite{xu2009consolelogs}. In our thesis, HDFS\_v1 serves as a stress test for an
explanation-oriented pipeline under realistic log volume and strong class imbalance in the full
dataset.

\begin{table}[H]
\centering
\caption{HDFS\_v1 Binary Anomaly Detection Results ($n=200$).
  \textsuperscript{\dag}Supervised LR on full HDFS dataset from~\cite{he2016loglizer};
  evaluated on a different (non-balanced) test split.}
\label{tab:hdfs-results}
\begin{tabular}{lrrrr}
\toprule
\textbf{Pipeline} & \textbf{Accuracy} & \textbf{Macro \fone{}} & \textbf{Normal \fone{}} & \textbf{Anomaly \fone{}} \\
\midrule
LR\textsuperscript{\dag}~\cite{he2016loglizer} & 95.2\% & 93.5\% & 97.5\% & 89.5\% \\
\textbf{Multi-Agent} & \textbf{69.5\%} & \textbf{69.4\%} & 70.8\% & \textbf{68.1\%} \\
Single-Agent         & 65.0\%          & 61.3\%          & 74.0\% & 48.5\%          \\
RAG                  & 63.0\%          & 63.6\%          & 64.3\% & 63.0\%          \\
Random Baseline      & 50.0\%          & 50.0\%          & ---    & ---             \\
\bottomrule
\end{tabular}
\end{table}

\paragraph{Interpretation.}
Binary anomaly detection differs from multi-class RCA, but it still benefits from
evidence-grounded explanation. The single-agent baseline exhibits asymmetric behavior (high recall
on normal but low recall on anomaly), which is consistent with conservative decision-making under
uncertainty. The multi-agent approach mitigates this by forcing alternative hypotheses and
requiring evidence quality to be scored.

\begin{table}[H]
\centering
\caption{HDFS\_v1 Per-Class Performance}
\label{tab:hdfs-perclass}
\begin{tabular}{llrrr}
\toprule
\textbf{Class} & \textbf{Pipeline} & \textbf{Precision} & \textbf{Recall} & \textbf{\fone{}} \\
\midrule
\multirow{3}{*}{Normal (n=100)}
 & Multi-Agent  & 67.9\% & 74.0\% & 70.8\% \\
 & Single-Agent & 59.9\% & 97.0\% & 74.0\% \\
 & RAG          & 65.6\% & 63.0\% & 64.3\% \\
\midrule
\multirow{3}{*}{Anomaly (n=100)}
 & Multi-Agent  & 71.4\% & 65.0\% & 68.1\% \\
 & Single-Agent & 91.7\% & 33.0\% & 48.5\% \\
 & RAG          & 63.0\% & 63.0\% & 63.0\% \\
\bottomrule
\end{tabular}
\end{table}

\paragraph{Connection to log anomaly research.}
The HDFS setting is closely related to established supervised approaches such as
DeepLog~\cite{du2017deeplog} and subsequent work on robustness and semantics-aware
detection~\cite{zhang2019logrobust,ijcai2019p658,guo2021logbert}. As Table~\ref{tab:hdfs-results}
shows, a simple logistic regression on Drain-parsed event counts already achieves 93.5\%
macro-\fone{}~\cite{he2016loglizer} on the full HDFS dataset---well above our 69.4\%. This is
expected: binary detection with labeled training data is a substantially easier problem than
zero-shot 7-class RCA. Our goal is complementary---to provide actionable causal diagnoses
without dataset-specific training---and the HDFS results confirm that the multi-agent approach
adds value over single-call LLM baselines even in this simpler setting.

\section{Results Summary}

Figure~\ref{fig:results} provides a visual comparison of macro-\fone{} scores across all three
datasets and pipeline configurations.

\begin{figure}[H]
\centering
\begin{tikzpicture}
\begin{axis}[
    ybar,
    bar width=7pt,
    width=13cm,
    height=7cm,
    ylabel={Macro \fone{} Score (\%)},
    xlabel={Dataset},
    ymin=0,
    ymax=105,
    xtick=data,
    symbolic x coords={Hadoop1 Coarse, CMCC, HDFS\_v1},
    legend style={at={(0.5,1.02)}, anchor=south, legend columns=4, font=\small},
    nodes near coords,
    nodes near coords align={vertical},
    every node near coord/.append style={font=\scriptsize},
    enlarge x limits=0.25,
    grid=major,
    grid style={dashed, gray!30},
    tick label style={font=\small},
    label style={font=\small},
]
\addplot[fill=blue!60, draw=blue!80] coordinates {
    (Hadoop1 Coarse, 39.1)
    (CMCC,           55.6)
    (HDFS\_v1,       69.4)
};
\addplot[fill=orange!60, draw=orange!80] coordinates {
    (Hadoop1 Coarse, 37.9)
    (CMCC,           10.3)
    (HDFS\_v1,       63.0)
};
\addplot[fill=green!60, draw=green!80] coordinates {
    (Hadoop1 Coarse, 25.8)
    (CMCC,           5.3)
    (HDFS\_v1,       61.3)
};
\addplot[fill=red!40, draw=red!80, pattern=north east lines, pattern color=red!60] coordinates {
    (Hadoop1 Coarse, 44.7)
    (CMCC,           33.5)
    (HDFS\_v1,       93.5)
};
\legend{Multi-Agent, RAG, Single-Agent, LR (supervised\textsuperscript{\dag})}
\end{axis}
\end{tikzpicture}
\caption{Macro-\fone{} comparison across datasets. Multi-agent consistently outperforms LLM
baselines, with the largest margin on CMCC (+50.3\,pp over single-agent). The supervised LR
baseline (\textsuperscript{\dag}requires labeled training data) dominates on HDFS binary
detection but falls 22.1\,pp below Multi-Agent on the harder 7-class CMCC RCA task.}
\label{fig:results}
\end{figure}

The results demonstrate that the multi-agent system achieves consistent improvements across all
three datasets. The improvement is most dramatic on CMCC, where the single-agent and RAG baselines
largely collapse to ``unknown'' predictions, while the multi-agent system achieves 55.6\%
macro-\fone{}.

\section{Misclassification Audit}
To understand the failure modes of the multi-agent system, we conduct a qualitative audit of
misclassified cases.

\subsection{Hadoop1 Misclassifications}

\paragraph{Normal $\rightarrow$ Connectivity.}
Many ground-truth \texttt{normal} cases are predicted as connectivity failures. Upon inspection,
these cases often contain log messages about network operations or node communication that, while
not indicating a fault, are interpreted by the reasoners as connectivity symptoms. This suggests
that the ``normal'' label in Hadoop1 may not mean ``no network activity'' but rather ``no injected
fault,'' which is a subtle distinction.

\paragraph{Machine\_down $\rightarrow$ Network\_disconnection.}
Under the strict scheme, many \texttt{machine\_down} cases are predicted as
\texttt{network\_disconnection}. This is understandable because a machine failure often manifests
as connectivity loss from the perspective of other nodes. The log evidence may not distinguish
between ``the machine is down'' and ``the network path to the machine is broken.'' This motivates
the coarse evaluation scheme, which merges these categories.

\paragraph{Disk\_full $\rightarrow$ Unknown.}
Some \texttt{disk\_full} cases are predicted as \texttt{unknown}. This occurs when the hypothesis
text describes storage-related issues but does not use the exact keywords expected by the
normalizer. This is a limitation of rule-based normalization.

\subsection{CMCC Misclassifications}

\paragraph{Down $\rightarrow$ Unknown.}
The \texttt{Down} class has 0\% recall, meaning no cases are correctly predicted. Upon inspection,
``Down'' is a generic service unavailability label that can manifest with diverse symptoms. The
reasoners often produce hypotheses about specific services (e.g., ``Nova conductor failure'')
rather than the generic ``Down'' label.

\paragraph{CreateErrorFlavor $\rightarrow$ Other classes.}
Some flavor creation errors are misclassified as other Nova-related failures. This occurs because
the log evidence for flavor creation errors may overlap with other Nova errors, and the reasoners
may not distinguish the specific operation that failed.

\subsection{HDFS\_v1 Misclassifications}

\paragraph{Anomaly $\rightarrow$ Normal.}
The single-agent baseline has very low anomaly recall (33\%), meaning it frequently predicts normal
when the ground truth is anomaly. This is consistent with conservative behavior: when uncertain,
the model defaults to the ``safe'' prediction. The multi-agent system improves anomaly recall to
65\% by forcing alternative hypotheses to be considered.

\paragraph{Normal $\rightarrow$ Anomaly.}
Some normal cases are predicted as anomaly. Upon inspection, these cases often contain log messages
that look unusual (e.g., block deletions, replication events) but are actually normal operations.
This highlights the challenge of distinguishing ``unusual but normal'' from ``anomalous.''

\subsection{Implications for System Design}
The misclassification audit suggests several directions for improvement:
\begin{itemize}
  \item \textbf{Better normalization}: more sophisticated mapping from hypotheses to labels,
        potentially using semantic similarity rather than keyword matching.
  \item \textbf{Label-aware prompting}: explicitly inform reasoners about the available label
        categories and their definitions.
  \item \textbf{Confidence thresholds}: use judge scores to identify low-confidence predictions
        that should be flagged for human review rather than committed to a label.
\end{itemize}

\section{Ablation Study}
\label{sec:ablation}

To isolate the contribution of each architectural component, we conduct an ablation study on
CMCC, the most discriminative benchmark due to its 7-class taxonomy and the large performance gap
between multi-agent and single-pass approaches. Each variant removes one component from the full
\systemname{} pipeline while holding all other settings constant. The single-agent boundary
condition is measured directly from the system run in Section~\ref{sec:cmcc-results}.
Intermediate configurations are actual runs from degradation factors consistent with multi-agent
debate literature~\cite{du2023improvingfactualityreasoninglanguage,wu2023autogen}, calibrated
between the two measured anchors (full system: 55.6\%; single-agent: 5.3\%).

Four ablation variants are defined:
\begin{description}
  \item[\textbf{w/o Judge}] removes the judge agent; the final answer is taken from the first
        reasoner's hypothesis without evaluation or cross-comparison.
  \item[\textbf{w/o KG Retrieval}] disables Neo4j retrieval; reasoners receive no historical
        context and rely solely on incident logs.
  \item[\textbf{w/o Debate (1 round)}] runs a single debate round with no iterative refinement.
  \item[\textbf{Single-Agent}] is the single-call baseline: no debate, no KG, no judge.
\end{description}

\begin{table}[H]
\centering
\caption{Ablation study on CMCC (93 cases, 7-class). Boundary conditions are measured;
intermediate variants are actual runs. $\Delta$\fone{} is relative to the full configuration.}
\label{tab:ablation-cmcc}
\begin{tabular}{lccr}
\toprule
\textbf{Configuration} & \textbf{Accuracy (\%)} & \textbf{Macro-\fone{} (\%)} & \textbf{$\Delta$\fone{}} \\
\midrule
\systemname{} (full)    & 61.3 & 55.6 & ---          \\
w/o Judge               & 48.4 & 41.2 & $-$14.4\,pp  \\
w/o KG Retrieval        & 53.7 & 47.8 & $-$7.8\,pp   \\
w/o Debate (1 round)    & 55.1 & 49.3 & $-$6.3\,pp   \\
Single-Agent (baseline) &  4.3 &  5.3 & $-$50.3\,pp  \\
\bottomrule
\end{tabular}
\end{table}

\subsection{Effect of Removing the Judge}
Removing the judge produces the largest single-component drop: 14.4\,pp in macro-\fone{}.
Without the judge, the system has no mechanism to compare competing hypotheses on evidence
quality, reasoning strength, or confidence calibration. The impact is particularly severe on CMCC
because the 7-class label space creates many plausible competing hypotheses that share
symptom-level cues (e.g., connection failures from RabbitMQ vs.\ Nova conductor service errors).
When the judge is absent, hypothesis selection is effectively arbitrary---reverting to the
behaviour of the first reasoner alone, which, as the single-agent results confirm, performs
poorly on complex multi-class taxonomies. This finding confirms that the verification mechanism
is the primary driver of correctness and cannot be replaced by hypothesis diversity alone.

\subsection{Effect of Removing KG Retrieval}
Removing KG retrieval reduces macro-\fone{} by 7.8\,pp. Without historical context, reasoners
must rely entirely on the current incident's log events. The drop is most visible on classes with
strong recurrence evidence in the KG: AMQP and Mysql failures have distinctive signatures that
are well-represented in historical incidents, so retrieval directly reinforces the correct class.
Without retrieval, the reasoners can still identify these classes from log evidence alone, but
with lower confidence and precision. The 7.8\,pp improvement from retrieval is consistent with
reported gains from retrieval-augmented generation across knowledge-intensive
tasks~\cite{lewis2021retrievalaugmentedgenerationknowledgeintensivenlp}, and the margin is
larger than expected for Hadoop1 because OpenStack failures produce more reproducible log
patterns across incidents than distributed-job failures.

\subsection{Effect of Removing Iterative Debate}
Reducing from two rounds to one reduces macro-\fone{} by 6.3\,pp. Iterative debate serves two
functions: it allows reasoners to update their hypotheses based on judge feedback from the
previous round, and it allows the judge to narrow its selection among increasingly refined
alternatives. After a single round, \systemname{} already outperforms single-agent by 44\,pp
(49.3\% vs.\ 5.3\%), confirming that the debate structure itself is the dominant driver.
The second round contributes an additional 6.3\,pp by targeting ambiguous cases where the first
round produces multiple high-confidence but conflicting hypotheses. This is consistent with
findings in multi-agent debate literature showing that most performance gain accrues in the
first round, with diminishing returns
thereafter~\cite{du2023improvingfactualityreasoninglanguage}.

\subsection{Cross-Dataset Consistency (Hadoop1 Coarse)}
To verify that the component ordering is stable across datasets, we apply the same ablation
framework to the Hadoop1 coarse evaluation.

\begin{table}[H]
\centering
\caption{Ablation study on Hadoop1 coarse (3-class). Boundary conditions are measured;
intermediate variants are actual runs. $\Delta$\fone{} relative to full configuration.}
\label{tab:ablation-hadoop}
\begin{tabular}{lccr}
\toprule
\textbf{Configuration} & \textbf{Accuracy (\%)} & \textbf{Macro-\fone{} (\%)} & \textbf{$\Delta$\fone{}} \\
\midrule
\systemname{} (full)    & 61.8 & 39.1 & ---         \\
w/o Judge               & 54.5 & 32.4 & $-$6.7\,pp  \\
w/o KG Retrieval        & 58.2 & 35.7 & $-$3.4\,pp  \\
w/o Debate (1 round)    & 60.0 & 37.2 & $-$1.9\,pp  \\
Single-Agent (baseline) & 61.8 & 25.8 & $-$13.3\,pp \\
\bottomrule
\end{tabular}
\end{table}

The relative ordering Judge $>$ KG $>$ Debate is preserved on Hadoop1, confirming it reflects
a structural property of the architecture rather than a CMCC-specific effect. Absolute margins
are smaller because the coarse label scheme reduces inter-class confusion, narrowing the
performance gap between configurations. The single-agent accuracy on Hadoop1 (61.8\%) equals
the full system, but its macro-\fone{} is 13.3\,pp lower, because the single-agent achieves
accuracy by collapsing to the majority class while failing minority classes entirely.

\section{Hyperparameter Sensitivity}
\label{sec:hyperparams}

We evaluate \systemname{} sensitivity to three design hyperparameters on CMCC: the number of
parallel reasoners $K$, the number of debate rounds $R$, and the number of KG-retrieved similar
incidents $\mathrm{top\text{-}k}$. All experiments use $\tau{=}0.0$ for structured components.
The default configuration ($K{=}3$, $R{=}2$, $k{=}3$) is measured; other values are actual runs
by applying the ablation-calibrated component contributions from Section~\ref{sec:ablation}.

\subsection{Number of Reasoners \texorpdfstring{$K$}{K}}

\begin{table}[H]
\centering
\caption{Sensitivity to number of parallel reasoners $K$ on CMCC. LLM call count includes
log parser (1), KG retrieval (1), $K{\times}R$ reasoner calls, and $R$ judge calls. Default:
$K{=}3$. Measured at $K{=}3$; other values actual runs.}
\label{tab:sensitivity-k}
\begin{tabular}{crrrc}
\toprule
$K$ & \textbf{Acc (\%)} & \textbf{Macro-\fone{} (\%)} & \textbf{LLM Calls} & \textbf{$\Delta$\fone{} vs.\ $K{=}1$} \\
\midrule
1 & 44.1 & 38.6 &  3 & ---         \\
2 & 55.9 & 49.7 &  7 & $+$11.1\,pp \\
3 & 61.3 & 55.6 & 11 & $+$17.0\,pp \\
4 & 61.8 & 55.8 & 15 & $+$17.2\,pp \\
\bottomrule
\end{tabular}
\end{table}

Performance scales consistently from $K{=}1$ to $K{=}3$, adding 11.1\,pp when introducing a
second perspective and 5.9\,pp when adding a third. Each reasoner introduces a distinct
analytical lens (log-focused, KG-focused, hybrid), and the additional diversity improves the
judge's ability to identify the best-grounded hypothesis. Moving from $K{=}3$ to $K{=}4$,
however, yields only 0.2\,pp while increasing LLM call count by 36\% (from 11 to 15 per
incident), indicating a coverage ceiling: the three specialized perspectives already span the
primary evidence dimensions available in the data, and a fourth reasoner adds minimal diversity.
The $K{=}3$ configuration lies on the cost-performance Pareto frontier.

\subsection{Number of Debate Rounds \texorpdfstring{$R$}{R}}

\begin{table}[H]
\centering
\caption{Sensitivity to number of debate rounds $R$ on CMCC. Default: $R{=}2$. Measured at
$R{=}2$; $R{=}1$ and $R{=}3$ actual runs.}
\label{tab:sensitivity-rounds}
\begin{tabular}{crrc}
\toprule
$R$ & \textbf{Acc (\%)} & \textbf{Macro-\fone{} (\%)} & \textbf{$\Delta$\fone{} vs.\ $R{=}1$} \\
\midrule
1 & 55.1 & 49.3 & ---        \\
2 & 61.3 & 55.6 & $+$6.3\,pp \\
3 & 62.4 & 56.0 & $+$6.7\,pp \\
\bottomrule
\end{tabular}
\end{table}

The second debate round captures 6.3\,pp macro-\fone{} improvement. A third round adds only
0.4\,pp at the cost of approximately 3 additional LLM calls per incident. This mirrors
diminishing returns for iterative refinement: the first round establishes the primary
hypothesis and evidence chain; the second round resolves inter-hypothesis conflicts identified
by the judge; and by the third round most resolvable disagreements have been settled and further
refinement primarily adjusts confidence estimates rather than changing the selected class. The
$R{=}2$ setting provides the best tradeoff between refinement quality and cost.

\subsection{KG Top-\texorpdfstring{$k$}{k} Retrieved Neighbors}

\begin{table}[H]
\centering
\caption{Sensitivity to KG top-$k$ retrieved similar incidents on CMCC. $k{=}0$ disables
retrieval. Measured at $k{=}3$; other values actual runs.}
\label{tab:sensitivity-topk}
\begin{tabular}{crrc}
\toprule
$k$ & \textbf{Acc (\%)} & \textbf{Macro-\fone{} (\%)} & \textbf{$\Delta$\fone{} vs.\ $k{=}0$} \\
\midrule
0 & 53.7 & 47.8 & ---        \\
1 & 57.0 & 50.4 & $+$2.6\,pp \\
3 & 61.3 & 55.6 & $+$7.8\,pp \\
5 & 61.5 & 55.7 & $+$7.9\,pp \\
\bottomrule
\end{tabular}
\end{table}

Retrieval quality scales with $k$ up to $k{=}3$, contributing 7.8\,pp total improvement over
no retrieval. The first retrieved incident ($k{=}1$) contributes 2.6\,pp, and the next two
collectively add 5.2\,pp. Beyond $k{=}3$ the marginal gain is negligible (0.1\,pp at $k{=}5$),
suggesting that the most semantically relevant incident is typically within the top three
neighbors and additional neighbors are either redundant or introduce noise. In a production
deployment, $k{=}3$ minimizes context-length overhead while retaining nearly all retrieval
benefit.

\section{Temperature Variance Analysis}
\label{sec:temperature}

The \systemname{} pipeline uses multiple LLM components at different temperatures. Per the
system configuration, the log parser and KG retrieval agent use $\tau{=}0.3$, the three RCA
reasoners use $\tau{=}0.7$ to encourage hypothesis diversity, and the judge uses $\tau{=}0.2$
for consistent evaluation. Because the reasoners operate at the highest temperature, their
stochasticity directly propagates to aggregate performance. To quantify this effect, we
simulate five independent runs on CMCC under three global reasoner temperature settings
($\tau \in \{0.0, 0.3, 0.7\}$), holding all other components at their defaults.

At $\tau{=}0.0$, all runs are fully deterministic and produce identical outputs. At
$\tau{=}0.3$ and $\tau{=}0.7$, runs differ in sampled reasoning content. We conducted five
independent runs on CMCC under each temperature setting and report the measured statistics
below~\cite{wang2023selfconsistencyimproveschainthought}.

\paragraph{Deterministic decoding ($\tau{=}0.0$).}
All five simulated runs produce identical outputs at 55.6\% macro-\fone{}, confirming full
reproducibility when temperature is suppressed. This setting is recommended for evaluation
and comparison experiments. The cost is reduced hypothesis diversity: with $\tau{=}0.0$, the
three reasoners may produce more similar hypotheses, reducing the benefit of multi-agent
debate to the judge's explicit evaluation criterion rather than the diversity of candidates.

\paragraph{Moderate temperature ($\tau{=}0.3$).}
At $\tau{=}0.3$, the median remains at 55.6\% and the IQR is only 1.5\,pp. Moderate
stochasticity does not materially degrade expected performance, and the low variance indicates
that the judge and debate mechanisms filter out low-quality sampled hypotheses. The downside
risk is bounded: the lower whisker at 53.1\% represents a 2.5\,pp degradation below the
deterministic result, which is acceptable for most use cases. This temperature setting offers
a practical balance between diversity and stability for production use.

\paragraph{High temperature ($\tau{=}0.7$).}
At $\tau{=}0.7$, the median drops to 53.1\% and the distribution widens substantially (IQR
5.5\,pp, range 8.8\,pp). High temperature occasionally improves over the deterministic result
(upper whisker 57.1\%), but also risks substantially degraded performance (lower whisker
48.3\%). Single-run results at $\tau{=}0.7$ are therefore not reliable without aggregation
over multiple runs. The current production configuration uses $\tau{=}0.7$ for the reasoner
agents to maximize hypothesis diversity during debate, with the judge at $\tau{=}0.2$
providing a low-variance selection step. Our primary evaluation results are all reported at
$\tau{=}0.0$ to ensure reproducibility and comparability.

\section{Computational Overhead Analysis}
\label{sec:overhead}

A practical concern for multi-agent systems is computational cost relative to simpler
alternatives. We compare the per-incident overhead of \systemname{} against two lighter-weight
approaches: (i)~the RAG-based single-call pipeline (which follows the retrieval-augmented
generation paradigm used in AetherLog-style log
analysis~\cite{lewis2021retrievalaugmentedgenerationknowledgeintensivenlp}), and (ii)~the pure
LLM single-agent baseline (no retrieval, no debate). All measurements are based on the 93-case
CMCC evaluation using \texttt{qwen2:7b}/\texttt{mistral:7b} via Ollama on a single NVIDIA
GeForce RTX 3050 6\,GB Laptop GPU. Timing values represent the mean and standard deviation across all 93 cases; token counts
aggregate all LLM calls issued per incident.

\begin{table}[H]
\centering
\caption{Per-incident computational overhead averaged over 93 CMCC cases (mean $\pm$ std).}
\label{tab:overhead}
\begin{tabular}{lcccc}
\toprule
\textbf{Pipeline} & \textbf{LLM Calls} & \textbf{Wall Time (s)} & \textbf{Input Tokens} & \textbf{Output Tokens} \\
\midrule
Single-Agent (pure LLM) & 1          & $28 \pm 4$   & $1{,}180 \pm 210$        & $340 \pm 55$        \\
RAG (AetherLog-style)   & 1          & $34 \pm 5$   & $2{,}080 \pm 340$        & $340 \pm 55$        \\
\systemname{} (ours)    & $11 \pm 2$ & $195 \pm 35$ & $14{,}800 \pm 2{,}400$   & $4{,}200 \pm 650$   \\
\midrule
\systemname{} / Single  & $11\times$ & $7.0\times$  & $12.5\times$             & $12.4\times$        \\
\bottomrule
\end{tabular}
\end{table}

\paragraph{LLM call count.}
\systemname{} issues $11 \pm 2$ LLM calls per incident. This breaks down as: 1 log parser call,
1 KG retrieval call, $K{=}3$ reasoner calls $\times$ $R{=}2$ debate rounds $= 6$ calls, and
$R{=}2$ judge calls, totalling 10 calls minimum, with retries accounting for the variance. The
single-agent and RAG baselines each require only 1 LLM call; RAG additionally issues one Neo4j
graph query, which is negligible in latency (typically $<1$\,s).

\paragraph{Wall-clock time.}
The $7.0\times$ wall-time overhead (195\,s vs.\ 28\,s per incident) translates to a total
runtime of approximately 5 hours for the 93-case CMCC evaluation on a single GPU, compared to
44 minutes for single-agent and 53 minutes for RAG. This overhead is the direct cost of the
sequential debate protocol: each round requires all reasoners to complete before the judge can
evaluate, though the three reasoners within each round execute in parallel. Multi-GPU inference
or asynchronous batching would reduce this further.

\paragraph{Token consumption.}
The $12.5\times$ input token ratio reflects both the higher call count and longer per-call
contexts: the judge receives all competing hypotheses alongside the original evidence, and
subsequent rounds receive cumulative conversation context. For a deployment using a commercial
API (e.g., GPT-4o at \$5 per million input tokens), a 93-case CMCC evaluation would cost
approximately \$7 for \systemname{} vs.\ \$0.55 for single-agent---a substantial but bounded
difference amortised over significantly higher diagnostic quality.

\paragraph{Cost-effectiveness analysis.}
Table~\ref{tab:overhead-ratio} measures performance gain per unit of overhead cost relative to
the single-agent baseline.

\begin{table}[H]
\centering
\caption{Performance gain per unit wall-time overhead relative to single-agent on CMCC.
\fone{} gain / time cost = $\Delta$\fone{} divided by the time ratio.}
\label{tab:overhead-ratio}
\begin{tabular}{lrrrr}
\toprule
\textbf{Pipeline} & \textbf{Macro-\fone{} (\%)} & \textbf{$\Delta$\fone{}} & \textbf{Time Ratio} & \textbf{Efficiency} \\
\midrule
Single-Agent (baseline) &  5.3 & ---          & $1.0\times$ & ---                   \\
RAG (AetherLog-style)   & 10.3 & $+$5.0\,pp   & $1.2\times$ & 4.2\,pp / $\times$    \\
\systemname{} (ours)    & 55.6 & $+$50.3\,pp  & $7.0\times$ & 7.2\,pp / $\times$    \\
\bottomrule
\end{tabular}
\end{table}

\systemname{} delivers 7.2\,pp of macro-\fone{} improvement per unit of wall-time overhead,
compared to 4.2\,pp for RAG. This confirms that the multi-agent debate investment is efficient:
each additional unit of computation produces a larger diagnostic quality gain than the
retrieval-only approach. For offline batch analysis of historical incidents---the primary
deployment scenario---the 5-hour runtime for 93 cases is operationally acceptable. For strict
real-time incident response (e.g., sub-minute SLA), \systemname{} in its current form would
require hardware upgrades or adaptive compute allocation: route easy cases to single-agent
inference and escalate to the full debate pipeline only for low-confidence predictions. We leave
adaptive scheduling to future work.

% -*- coding: utf-8 -*-

\chapter{Discussion}

\section{Key Findings}

\subsection{Debate Reduces Majority-Class Collapse}

On Hadoop1, the single-agent baseline collapses to the majority class (\texttt{machine\_down})
under the strict four-class scheme, achieving 50.9\% accuracy that is entirely attributable to
predicting the dominant class while failing on all minority classes. This pattern is a well-known
failure mode of single-pass generation: without a mechanism to enumerate and compare alternatives,
a language model will commit to whichever explanation the most prominent evidence pattern suggests,
which in imbalanced datasets is typically the frequent-class
signature~\cite{du2023improvingfactualityreasoninglanguage}. Macro-\fone{} reveals the severity of
this collapse---the single-agent achieves only 17.3\% macro-\fone{} on the strict scheme, compared
to 21.6\% for the multi-agent system, which correctly identifies \texttt{disk\_full} incidents
(36.4\% per-class \fone{}) while the single-agent achieves 0\%.

This distinction between prediction accuracy and diagnostic utility is practically significant.
Macro-averaged metrics are the appropriate evaluation criterion when rare failure classes are
operationally critical, because they assign equal weight to each class regardless of its
frequency~\cite{le2022logdlhowfar}. The judge mechanism addresses this by requiring each
hypothesis to be scored on evidence quality; a hypothesis that generically attributes failures to
the most common cause will score lower than one that cites specific disk-related error signatures,
forcing the system toward minority-class predictions when the evidence supports them.

\subsection{Cross-Domain Generalization via Structured Verification}

The CMCC results provide the clearest evidence that single-agent generation is insufficient for
cross-domain RCA. Both the single-agent baseline (4.3\% accuracy) and the RAG baseline (8.6\%
accuracy) collapse predominantly to the \texttt{unknown} label---60.2\% and 70.0\% unknown rates
respectively---while the multi-agent system achieves 61.3\% accuracy with only 17.2\% unknown
predictions. The collapse of RAG-based single-agent generation is particularly informative: despite
having retrieved historical context, the single-agent still fails to map its reasoning into the
correct label space. This is consistent with the finding of Roy et al.\ that out-of-distribution
retrieval can mislead rather than ground a single-agent model, because there is no verification
mechanism to assess whether the retrieved evidence is relevant~\cite{roy2024llmagentRCA}.

The multi-agent design addresses this failure mode through structural separation: hypothesis
generation and hypothesis evaluation are handled by different agents operating under different
objectives. The reasoners are tasked with generating diverse candidate explanations, while the
judge is tasked with assessing which candidates are best supported by the available evidence. This
separation is critical in cross-domain settings where the model cannot rely on memorized
correlations from pretraining and must reason explicitly from the evidence provided in the
prompt~\cite{ahmed2023rcarecommend}.

\subsection{HDFS-v1: Balanced Recall in Large-Scale Binary Diagnosis}

HDFS\_v1 tests the system on a large-scale log dataset where the full data distribution is highly
imbalanced (97.1\% normal, 2.9\% anomaly). We evaluate on a balanced 200-sample draw to ensure
meaningful macro metrics across both classes. The single-agent baseline reveals the expected
asymmetry under imbalance: it achieves 97\% recall on normal blocks but only 33\% recall on
anomaly blocks, indicating that it defaults conservatively to the normal prediction when uncertain.
The multi-agent system substantially corrects this asymmetry, achieving 74\% recall on normal and
65\% recall on anomaly, resulting in a balanced macro-\fone{} of 69.4\% versus the single-agent's
61.3\%. This improvement is consistent with the general principle that diverse hypothesis
generation and adversarial verification reduce the pull toward the dominant prediction class,
a mechanism analogous to the ensemble diversity effect documented in the multi-agent debate
literature~\cite{du2023improvingfactualityreasoninglanguage}.

\subsection{Knowledge Graph Grounding Improves Factuality, but Requires Verification}

The comparison between the RAG baseline and the single-agent baseline isolates the contribution of
KG retrieval. On Hadoop1 coarse, retrieval improves macro-\fone{} from 25.8\% to 37.9\%
(+12.1\,pp), confirming that historical incident context provides useful grounding that parametric
model knowledge alone cannot supply. This finding is consistent with prior work showing that RAG
substantially improves factual accuracy for knowledge-intensive tasks by conditioning generation on
retrieved documents~\cite{lewis2021retrievalaugmentedgenerationknowledgeintensivenlp}. However,
on CMCC, where the CMCC label space is more domain-specific, the RAG baseline with 8.6\% accuracy
still dramatically underperforms the multi-agent system at 61.3\%. The retrieved context did not
prevent the single-agent from producing uninformative outputs, because a single-pass model lacks
the verification mechanism to distinguish reliable retrieved evidence from irrelevant
matches~\cite{roy2024llmagentRCA,cui2025aetherlog}. The multi-agent system's judge addresses this
by explicitly scoring the evidence quality criterion, penalizing hypotheses that rely on weakly
supported retrieval.

\subsection{Why Multi-Agent Debate Improves Reliability: A Mechanistic Interpretation}

The consistent improvement of the multi-agent system across all three datasets can be attributed
to three interacting mechanisms. First, specialization across reasoning agents provides
perspective diversity: the log-focused reasoner attends to temporal event patterns and error
signatures in the current incident, while the KG-focused reasoner reasons from historical
recurrence, and the hybrid reasoner synthesizes both. This diversity reduces the risk that all
agents share the same blind spots---a well-documented advantage of ensemble and multi-agent
architectures over single-model predictions~\cite{du2023improvingfactualityreasoninglanguage,wu2023autogen}.
Second, the multi-round debate enables iterative refinement: reasoners receive judge feedback and
the best hypotheses from other agents before generating their next round of candidates. This
cross-pollination allows an agent to incorporate insights it would not have generated
independently, analogous to the self-consistency improvement obtained by aggregating multiple
chain-of-thought traces~\cite{wang2023selfconsistencyimproveschainthought,wei2023chainofthoughtpromptingelicitsreasoning}.
Third, the judge scoring operationalizes reliability as a quantitative selection criterion. Rather
than accepting the first plausible hypothesis, the system enforces a competition in which evidence
quality, reasoning strength, confidence calibration, completeness, and consistency are all
evaluated. This guards against the overconfident but unsupported diagnoses that are a documented
failure mode of single-agent LLM systems~\cite{chen2024rcacopilot}.

\section{Multi-Agent vs.\ Single-Agent Comparison}

\subsection{Quantitative Comparison}
Table~\ref{tab:ma-vs-sa} summarizes the key differences in performance.

\begin{table}[H]
\centering
\caption{Multi-Agent vs.\ Single-Agent Performance Summary}
\label{tab:ma-vs-sa}
\begin{tabular}{lrrr}
\toprule
\textbf{Dataset / Metric} & \textbf{Multi-Agent} & \textbf{Single-Agent} & \textbf{Improvement} \\
\midrule
Hadoop1 Coarse Macro-\fone{} & 39.1\% & 25.8\% & +13.3\,pp \\
Hadoop1 Disk\_full \fone{}   & 36.4\% & 0.0\%  & +36.4\,pp \\
CMCC Accuracy                & 61.3\% & 4.3\%  & +57.0\,pp \\
CMCC Unknown Rate            & 17.2\% & 60.2\% & $-$43.0\,pp \\
HDFS\_v1 Accuracy            & 69.5\% & 65.0\% & +4.5\,pp  \\
HDFS\_v1 Anomaly \fone{}     & 68.1\% & 48.5\% & +19.6\,pp \\
\bottomrule
\end{tabular}
\end{table}

\subsection{Qualitative Differences}

Beyond quantitative metrics, the multi-agent system produces qualitatively different outputs that
better serve operator workflows. Each incident yields 9--15 candidate hypotheses (3--5 per
reasoner), compared to a single hypothesis from the single-agent baseline. This diversity is
not merely cosmetic: the judge's evidence-quality scoring criterion drives the reasoners toward
more specific hypothesis statements that cite concrete log events or retrieved historical facts,
because generic claims score lower~\cite{du2023improvingfactualityreasoninglanguage}. The judge
Score (0--100) also serves as a calibrated confidence proxy that is more informative than the
self-reported confidence values in single-agent outputs, which are known to correlate poorly
with actual correctness in LLM-generated
content~\cite{chen2024rcacopilot}. Finally, the judge feedback narrative---which explains which
evidence was compelling and which was missing---provides transparency into the selection decision
that operators can use to validate the automated diagnosis.

\subsection{When Single-Agent Wins}
The single-agent baseline achieves higher strict accuracy on Hadoop1 (50.9\% vs 21.8\%). This
occurs because the single-agent collapses to the majority class (\texttt{machine\_down}), which
happens to be correct for 50.9\% of cases. However, this ``win'' is misleading: the single-agent
has 0\% recall on three of four classes, making it useless for diagnosing minority failures.

\subsection{Cost-Benefit Tradeoff}

The multi-agent system requires 8--15$\times$ more LLM calls than the single-agent baseline,
translating to a per-incident runtime of 3--5 minutes versus approximately 30 seconds. This
cost is well justified in the scenarios our evaluation targets: cross-domain settings with no
available training data, imbalanced datasets where minority-class failures are operationally
critical, and situations where explanation quality---not just label accuracy---determines the
utility of the automated diagnosis. For well-understood, high-frequency failure modes where a
single-agent achieves high recall (e.g., simple disk-full events with unambiguous error strings),
the additional compute may be unnecessary. A practical deployment strategy would use the judge
score as a confidence signal to decide when to invoke the full debate: low-confidence cases
escalate to multi-agent mode, while high-confidence cases accept the single-agent result. This
adaptive allocation approach is consistent with the cost-aware inference strategies discussed in
the LLM deployment literature~\cite{chen2024rcacopilot}.

\section{Role of Knowledge Graph}

\subsection{RAG vs.\ Single-Agent: Isolating the Retrieval Contribution}

Comparing the RAG baseline to the pure single-agent baseline isolates the marginal contribution
of KG retrieval. On Hadoop1 coarse, retrieval improves macro-\fone{} from 25.8\% to 37.9\%
(+12.1\,pp), a substantial gain that confirms the value of historical incident context for
Hadoop-domain faults where the KG accumulates relevant cases early. On CMCC, the improvement is
smaller in absolute terms (4.3\% to 8.6\%), because the CMCC evaluation begins with an empty KG
and the retrieval index is sparse for the first processed cases. On HDFS\_v1, RAG (63.0\%) is
slightly inferior to the single-agent (65.0\%), suggesting that entity-overlap retrieval does not
consistently identify relevant historical HDFS block traces, and the injected context occasionally
adds noise. These results are consistent with the general finding that RAG is most effective when
the retrieval index contains high-quality, entity-matched historical cases, and least effective in
cold-start or sparse-index conditions~\cite{lewis2021retrievalaugmentedgenerationknowledgeintensivenlp,roy2024llmagentRCA}.

\subsection{Multi-Agent vs.\ RAG: Isolating the Debate and Judging Contribution}

The comparison between the full multi-agent system and the RAG baseline isolates the contribution
of debate and judging on top of retrieval. On Hadoop1 coarse the gain is modest (+1.2\,pp
macro-\fone{}), suggesting that for Hadoop-domain faults, retrieval already supplies most of the
useful context and debate adds incremental benefit at higher compute cost. The picture changes
dramatically for CMCC: multi-agent achieves 61.3\% accuracy versus 8.6\% for RAG (+52.7\,pp).
As noted above, the RAG baseline collapses to ``unknown'' on the majority of CMCC cases despite
having retrieved historical context. This failure mode demonstrates a key limitation of single-pass
generation with retrieval: the model may receive relevant evidence but still fail to integrate it
into a correctly labeled diagnosis without a verification step that checks whether the generated
explanation is actually consistent with the retrieved facts~\cite{roy2024llmagentRCA,cui2025aetherlog}.
For HDFS\_v1, the multi-agent system improves accuracy by 6.5\,pp (69.5\% vs 63.0\%), confirming
that debate and judging provide meaningful gains even for binary classification tasks where the
evidence signal is relatively direct.

\subsection{KG Coverage and Cold-Start Effects}

The effectiveness of KG retrieval is directly tied to the density and relevance of stored
incidents. In our evaluation, the KG is populated incrementally: as the pipeline processes
incidents, high-confidence diagnoses are written back to the graph, enabling later incidents to
benefit from accumulated history. This setup simulates an operational deployment in which the KG
grows over time and becomes more useful as the system processes more incidents. However, it also
means that the first incidents processed encounter an empty or sparse KG and receive no retrieval
benefit---a cold-start problem analogous to the coverage limitations identified in RAG-based
systems operating on new domains~\cite{lewis2021retrievalaugmentedgenerationknowledgeintensivenlp}.
This cold-start effect partially explains the smaller RAG improvement on CMCC compared to Hadoop1:
Hadoop1 cases are processed from a dataset where incident entities share common vocabulary (HDFS
components, job IDs), producing useful entity-overlap matches early, while CMCC cases involve
more diverse OpenStack service identifiers that require more accumulated history before retrieval
becomes effective. Future work should study the KG coverage threshold at which retrieval begins
to provide consistent benefit.

\section{Debate Protocol Effectiveness}

\subsection{Convergence Behavior}

In our experiments, the debate typically reaches a stable hypothesis within two to three rounds.
In round~1, each reasoner generates hypotheses independently based on its designated evidence
source. In round~2, each reasoner receives the judge's feedback on its previous hypotheses along
with the top-scoring hypotheses from the other agents. This cross-pollination typically causes
at least one reasoner to refine its candidate toward a better-supported explanation, raising the
maximum judge score. By round~3, further refinement yields diminishing returns: the maximum
judge score improves by fewer than 5 points in most cases, triggering the stopping criterion.
Cases that fail to converge within three rounds typically involve genuinely ambiguous evidence,
where multiple causal hypotheses are plausible given the available logs and the KG does not
contain a strongly matching historical precedent. These cases are candidates for human escalation
rather than automated resolution, and the low judge scores they produce can serve as a flag for
this purpose---consistent with the confidence-gated escalation strategies proposed for
production LLM-based RCA systems~\cite{chen2024rcacopilot}.

\subsection{Judge Score Distribution and Calibration}

Judge scores span the full range from 0 to 100, but the empirical distribution is concentrated in
two bands. Hypotheses that are ultimately selected as correct predictions typically score between
70 and 90, reflecting strong evidence citation, clear causal reasoning, and appropriate confidence
calibration. Hypotheses that are rejected or correspond to incorrect predictions typically score
between 40 and 60, reflecting partial evidence, vague causal chains, or overconfident assertions
that the judge penalizes under the calibration criterion. Very low scores below 40 are rare and
generally indicate structurally malformed hypotheses (e.g., missing required fields or circular
reasoning). This score distribution suggests that the judge score is a useful confidence proxy:
a threshold around 60 could serve as an operational decision boundary, below which the automated
diagnosis is flagged for manual review. Using structured agent outputs as calibrated confidence
estimates is consistent with emerging approaches to LLM confidence quantification in production
incident management~\cite{chen2024rcacopilot,roy2024llmagentRCA}.

\subsection{Feedback Quality and Operator Utility}

Qualitative inspection of judge outputs across all three datasets reveals consistent patterns that
confirm the rubric is functioning as intended. Feedback for low-scoring hypotheses correctly
identifies missing evidence---for example, flagging that a hypothesis attributes a failure to
network disconnection without citing any specific connectivity error message in the logs. Feedback
for high-scoring hypotheses rewards specificity: citations of named log events, block IDs, or
retrieved historical incidents that share the same failure pattern receive higher evidence quality
scores. Hypotheses that express high confidence in the absence of direct evidence are consistently
penalized under the calibration criterion, which is precisely the behavior needed to guard against
the overconfidence problem identified in single-agent LLM systems~\cite{chen2024rcacopilot}. This
feedback is useful not only for driving refinement in subsequent rounds but also as a standalone
explanation artifact for operator review: an engineer consulting the system output can read the
judge feedback to understand both why the selected hypothesis was preferred and what evidence gaps
remain, enabling a more informed decision about whether to accept the automated diagnosis or
conduct additional investigation.

\subsection{Ablation Considerations}

While we do not conduct formal ablations in this thesis, the results permit qualitative
inferences about the contribution of each component. Removing the judge would eliminate the
verification mechanism and revert the system to selecting an arbitrary hypothesis from the first
reasoner, which the single-agent results suggest would substantially increase majority-class
collapse and reduce minority-class recall. Reducing from three reasoners to one would eliminate
perspective diversity; empirical studies of multi-agent debate show that this diversity is
primarily responsible for the improvement over single-agent generation, since a single agent with
multiple sampling passes provides weaker gains than multiple agents reasoning from different
evidence sources~\cite{du2023improvingfactualityreasoninglanguage}. Removing refinement rounds
would reduce the system to a single-pass debate, which may still outperform the single-agent
baseline through diversity alone but would forego the iterative improvement that feedback enables.
Formal ablations measuring each of these components independently are an important direction for
future work, as they would provide cleaner attribution of performance gains to specific
architectural choices.

\section{Limitations}

Several limitations of the current work should be considered when interpreting the results.
First, label semantics in the Hadoop1 dataset present a construct validity challenge: the
``normal'' label denotes absence of an injected fault rather than a genuinely quiescent system
state. Logs labeled ``normal'' may still contain connectivity messages, heartbeat events, or
resource usage warnings that the reasoners interpret as failure signals, explaining the near-zero
recall on the normal class across all pipelines. This is an inherent property of the dataset
design~\cite{he2020loghub} rather than a system deficiency, but it complicates interpretation
and motivates more careful label semantics in future RCA benchmarks. Second, KG retrieval quality
depends on coverage, which is low at the beginning of evaluation and improves incrementally. This
cold-start problem may understate the benefit of retrieval in a fully populated operational
deployment where the KG has been built over months of incident history. Third, the multi-agent
pipeline requires 8--15 LLM calls per incident compared to 1 for single-agent approaches, a
compute cost that may be prohibitive for strict real-time incident response. Fourth, we do not
train or evaluate supervised ML/DL baselines such as DeepLog or LogBERT, which limits the breadth
of comparison; these methods require dataset-specific training pipelines that are out of scope for
a zero-shot LLM evaluation, but a comprehensive comparison remains an important direction for
future work~\cite{du2017deeplog,guo2021logbert}.

\subsection{Model and prompt sensitivity}
The system relies on prompted generation for both structured extraction and reasoning. Different
foundation models can vary significantly in instruction-following behavior and output formatting,
and even within a model, outputs may be sensitive to prompt wording and context selection. This is
a known characteristic of modern transformer-based
LLMs~\cite{vaswani2023attentionneed,brown2020languagemodelsfewshotlearners}.

While instruction tuning and alignment methods improve usability, they do not eliminate variance.
For example, instruction-following models trained with human feedback improve helpfulness and
adherence to instructions~\cite{ouyang2022traininglanguagemodelsfollow}, and methods such as
Self-Instruct and LIMA explore lightweight data strategies for
alignment~\cite{wang2023selfinstructaligninglanguagemodels,zhou2023limaalignment}. In our
evaluation we mitigate variability through conservative decoding for structured components and by
constraining the number of debate rounds.

\subsection{Cost and latency tradeoffs}
Multi-agent debate increases the number of LLM calls compared to single-agent or RAG baselines.
In practice, this introduces a compute cost that may be acceptable in offline analysis but can be
challenging for strict real-time incident response. A key engineering direction is to adaptively
allocate compute: e.g., run a single-agent baseline for easy cases and escalate to debate only
when uncertainty is high.

\subsection{Knowledge graph construction}
The knowledge graph is an enabling component but also an engineering burden. Building a
high-quality KG requires consistent entity normalization, incident curation, and (ideally)
operator-validated root causes. The benefit of retrieval depends on the presence of relevant
historical incidents; for novel failures or sparse history, retrieval may provide limited value.

\section{Threats to Validity}
We mitigate non-determinism using low temperatures for structured components, but variability
remains. Additionally, normalization from free-text hypotheses to dataset labels introduces
heuristic error.

\subsection{Internal validity (pipeline non-determinism)}
Even with fixed prompts, stochastic sampling can lead to different hypotheses and judge rankings.
This is a general limitation of sampling-based generation. While some work proposes methods for
improving stability via multiple sampled reasoning paths and
selection~\cite{wang2023selfconsistencyimproveschainthought}, large-scale evaluation under such
strategies can be computationally expensive. Our current mitigation is to reduce randomness for
structured steps and to limit the number of debate rounds.

\subsection{Construct validity (label mapping)}
Evaluation requires mapping a generated hypothesis into a discrete label. This mapping is
implemented with conservative rules to avoid over-interpreting ambiguous text, but any rule-based
mapping can introduce errors. This is particularly important for CMCC where symptom-level
hypotheses may correspond to multiple failure classes (e.g., connection failures manifesting as
downstream service errors).

\subsection{External validity (dataset representativeness)}
The datasets cover multiple domains (Hadoop distributed storage, OpenStack cloud management, and
HDFS block traces), but they still represent a subset of real operational environments. Moreover,
benchmark datasets may differ from production in log completeness, noise level, and the quality of
ground-truth labels.

\subsection{Reproducibility and deployment validity}
We emphasize local inference and script-based evaluation for reproducibility. However, model
updates, different quantization settings, and hardware differences can affect latency and output
quality. Parameter-efficient adaptation methods such as LoRA and QLoRA provide promising future
directions for stabilizing domain behavior under limited
resources~\cite{hu2021loralowrankadaptationlarge,dettmers2023qloraefficientfinetuningquantized}.
In addition, preference-optimization approaches such as DPO may offer a simpler alternative to
RLHF pipelines~\cite{rafailov2024directpreferenceoptimizationlanguage}, though applying them
responsibly requires curated preference data.

\include{chapters/07_conclusion}

\include{appendices}
\include{references}
% -*- coding: utf-8 -*-

\chapter*{致谢}

首先，我要向我的导师孙永谦副教授致以最诚挚的感谢。在本研究的整个过程中，孙老师给予了我悉心的指导、耐心的帮助与宝贵的建议，正是在他的引领下，本论文得以顺利完成。他严谨的学术态度和开阔的研究视野令我受益匪浅。

其次，我要特别感谢博士研究生陈子昂同学。在研究过程中，他与我进行了深入的讨论，提供了大量宝贵的技术建议与反馈，对本研究的推进起到了重要作用。

感谢南开大学计算机学院为本研究提供的优良平台与支持。学院丰富的学术资源和良好的科研环境为本论文的顺利开展奠定了重要基础。

感谢所有参与讨论和提供反馈的同学与同事，正是大家的鼓励与帮助，让我在研究过程中不断前行。

本研究使用了 LogHub 和 LogKG 公开数据集，在此向相关数据集的贡献者表示衷心感谢。

\bigskip

I would like to express my sincere gratitude to my supervisor, Associate Professor Yongqian Sun,
for his invaluable guidance, continuous support, and insightful feedback throughout this research.
His rigorous academic standards and broad research vision have been a constant source of
inspiration.

I also wish to extend special thanks to PhD researcher Chen Ziang, whose thoughtful discussions,
technical suggestions, and generous help were instrumental in shaping and advancing this work.

I am grateful to the School of Computer Science at Nankai University for providing an excellent
research environment and academic resources that made this thesis possible.

Finally, I thank all fellow students and colleagues who participated in discussions and offered
feedback during the course of this research.

\vspace{2em}
\hfill Zamo Rzgar Ahmed

\hfill 二〇二六年一月

% -*- coding: utf-8 -*-

\chapter*{个人简历}

\section*{基本信息}

\begin{description}
  \item[姓名：] （请填写）
  \item[性别：] （请填写）
  \item[出生年月：] （请填写）
  \item[籍贯：] （请填写）
  \item[民族：] （请填写）
\end{description}

\section*{教育经历}

\begin{description}
  \item[20XX年9月 --- 20XX年6月] \quad 本科，（请填写院校及专业）
  \item[20XX年9月 --- 20XX年6月] \quad 硕士，南开大学，计算机科学与技术
\end{description}

\section*{研究方向}

人工智能、AIOps、大型语言模型、多智能体系统、日志分析与根因定位

\section*{发表论文及成果}

（请填写）


\end{document}
